\RequirePackage[l2tabu,orthodox]{nag}

% TODO: decide if one-sided/two-sided
%\documentclass[headsepline,footsepline,footinclude=false,fontsize=11pt,paper=a4,listof=totoc,bibliography=totoc,BCOR=12mm,DIV=12]{scrbook} % two-sided
\documentclass[headsepline,footsepline,footinclude=false,oneside,fontsize=11pt,paper=a4,listof=totoc,bibliography=totoc]{scrbook} % one-sided

\PassOptionsToPackage{table,svgnames,dvipsnames}{xcolor}

\usepackage[utf8]{inputenc}
\usepackage[T1]{fontenc}
\usepackage[sc]{mathpazo}
\usepackage[ngerman,american]{babel}
\usepackage[autostyle]{csquotes}
\usepackage[%
  backend=biber,
  url=true,
  style=numeric,
  sorting=none,
  maxnames=4,
  minnames=3,
  maxbibnames=99,
  giveninits,
  uniquename=init]{biblatex}
\usepackage{graphicx}
\usepackage{scrhack} % necessary for listings package
\usepackage{listings}
\usepackage{lstautogobble}
\usepackage{tikz}
\usepackage{pgfplots}
\usepackage{pgfplotstable}
\usepackage{booktabs}
\usepackage[final]{microtype}
\usepackage{caption}
\usepackage{mathpartir}
\usepackage{mathrsfs}
\usepackage{mathtools}
\usepackage{amsthm}
\usepackage{amssymb}
\usepackage{stmaryrd}
\usepackage{turnstile}
%\usepackage[hidelinks]{hyperref} % hidelinks removes colored boxes around references and links
\usepackage{varwidth,pst-tree,realscripts}

% Isabelle specifics
\usepackage{isabelle,isabellesym}
% this should be the last package used
\usepackage{pdfsetup}

\bibliography{bibliography}


% Tree Figures
\psset{showbbox=false,treemode=D,linewidth=0.5pt,treesep=4ex,levelsep=1cm}
\newcommand{\LFTw}[2]{%
\Tr[ref=#1]{\psframebox[linestyle=none,framesep=4pt]{%
\begin{varwidth}{15ex}\center #2\end{varwidth}}}}
\newcommand{\LFTr}[2]{\Tr[ref=#1]{\psframebox[linestyle=none,framesep=4pt]{#2}}}

%\urlstyle{rm}
\isabellestyle{it}

\setkomafont{disposition}{\normalfont\bfseries} % use serif font for headings
\linespread{1.05} % adjust line spread for mathpazo font

% Add table of contents to PDF bookmarks
\BeforeTOCHead[toc]{{\cleardoublepage\pdfbookmark[0]{\contentsname}{toc}}}

% Define TUM corporate design colors
% Taken from http://portal.mytum.de/corporatedesign/index_print/vorlagen/index_farben
\definecolor{TUMBlue}{HTML}{0065BD}
\definecolor{TUMSecondaryBlue}{HTML}{005293}
\definecolor{TUMSecondaryBlue2}{HTML}{003359}
\definecolor{TUMBlack}{HTML}{000000}
\definecolor{TUMWhite}{HTML}{FFFFFF}
\definecolor{TUMDarkGray}{HTML}{333333}
\definecolor{TUMGray}{HTML}{808080}
\definecolor{TUMLightGray}{HTML}{CCCCC6}
\definecolor{TUMAccentGray}{HTML}{DAD7CB}
\definecolor{TUMAccentOrange}{HTML}{E37222}
\definecolor{TUMAccentGreen}{HTML}{A2AD00}
\definecolor{TUMAccentLightBlue}{HTML}{98C6EA}
\definecolor{TUMAccentBlue}{HTML}{64A0C8}

% Settings for pgfplots
\pgfplotsset{compat=newest}
\pgfplotsset{
  % For available color names, see http://www.latextemplates.com/svgnames-colors
  cycle list={TUMBlue\\TUMAccentOrange\\TUMAccentGreen\\TUMSecondaryBlue2\\TUMDarkGray\\},
}

% Settings for lstlistings
\lstset{%
  basicstyle=\ttfamily,
  columns=fullflexible,
  autogobble,
  keywordstyle=\bfseries\color{TUMBlue},
  stringstyle=\color{TUMAccentGreen}
}

% TODO: change thesis information
\newcommand*{\getUniversity}{Technische Universität München}
\newcommand*{\getFaculty}{Department of Informatics}
\newcommand*{\getTitle}{Formal Verification of an Earley Parser}
\newcommand*{\getTitleGer}{Formale Verifikation eines Earley Parsers}
\newcommand*{\getAuthor}{Martin Rau}
\newcommand*{\getDoctype}{Master's Thesis in Informatics}
\newcommand*{\getSupervisor}{Tobias Nipkow}
\newcommand*{\getAdvisor}{Tobias Nipkow}
\newcommand*{\getSubmissionDate}{15.06.2023}
\newcommand*{\getSubmissionLocation}{Munich}

\begin{document}

% Set page numbering to avoid "destination with the same identifier has been already used" warning for cover page.
% (see https://en.wikibooks.org/wiki/LaTeX/Hyperlinks#Problems_with_Links_and_Pages).
\pagenumbering{alph}

%%%%%%%%%%%%%%%%%%%%%%%%%%%%%%%%%%%%%%%%%%%%%%%%%%%%%%%%%%%%%%
%                                                            %
%                                                            %
%                         COVER                              %
%                                                            %
%                                                            %
%%%%%%%%%%%%%%%%%%%%%%%%%%%%%%%%%%%%%%%%%%%%%%%%%%%%%%%%%%%%%%

%\input{pages/cover}
\begin{titlepage}
  % HACK for two-sided documents: ignore binding correction for cover page.
  % Adapted from Markus Kohm's KOMA-Script titlepage=firstiscover handling.
  % See http://mirrors.ctan.org/macros/latex/contrib/koma-script/scrkernel-title.dtx,
  % \maketitle macro.
  \oddsidemargin=\evensidemargin\relax
  \textwidth=\dimexpr\paperwidth-2\evensidemargin-2in\relax
  \hsize=\textwidth\relax

  \centering

  \IfFileExists{tum.pdf}{%
    \includegraphics[height=20mm]{tum.pdf}
  }{%
    \vspace*{20mm}
  }

  \vspace{5mm}
  {\huge\MakeUppercase{\getFaculty{}}}\\

  \vspace{5mm}
  {\large\MakeUppercase{\getUniversity{}}}\\

  \vspace{20mm}
  {\Large \getDoctype{}}

  \vspace{15mm}
  {\huge\bfseries \getTitle{}}

  \vspace{15mm}
  {\LARGE \getAuthor{}}

  \IfFileExists{faculty.pdf}{%
    \vfill{}
    \includegraphics[height=20mm]{faculty.pdf}
  }{}
\end{titlepage}

\frontmatter{}

%%%%%%%%%%%%%%%%%%%%%%%%%%%%%%%%%%%%%%%%%%%%%%%%%%%%%%%%%%%%%%
%                                                            %
%                                                            %
%                         TITLE                              %
%                                                            %
%                                                            %
%%%%%%%%%%%%%%%%%%%%%%%%%%%%%%%%%%%%%%%%%%%%%%%%%%%%%%%%%%%%%%

\begin{titlepage}
  \centering

  \IfFileExists{tum.pdf}{%
    \includegraphics[height=20mm]{tum.pdf}
  }{%
    \vspace*{20mm}
  }

  \vspace{5mm}
  {\huge\MakeUppercase{\getFaculty{}}}\\

  \vspace{5mm}
  {\large\MakeUppercase{\getUniversity{}}}\\

  \vspace{20mm}
  {\Large \getDoctype{}}

  \vspace{15mm}
  {\huge\bfseries \getTitle{} \par}

  \vspace{10mm}
  {\huge\bfseries \foreignlanguage{ngerman}{\getTitleGer{}} \par}

  \vspace{15mm}
  \begin{tabular}{l l}
    Author:          & \getAuthor{} \\
    Supervisor:      & \getSupervisor{} \\
    Advisor:         & \getAdvisor{} \\
    Submission Date: & \getSubmissionDate{} \\
  \end{tabular}

  \IfFileExists{faculty.pdf}{%
    \vfill{}
    \includegraphics[height=20mm]{faculty.pdf}
  }{}
\end{titlepage}

%%%%%%%%%%%%%%%%%%%%%%%%%%%%%%%%%%%%%%%%%%%%%%%%%%%%%%%%%%%%%%
%                                                            %
%                                                            %
%                       DISCLAIMER                           %
%                                                            %
%                                                            %
%%%%%%%%%%%%%%%%%%%%%%%%%%%%%%%%%%%%%%%%%%%%%%%%%%%%%%%%%%%%%%

\thispagestyle{empty}
\vspace*{0.8\textheight}
\noindent
I confirm that this \MakeLowercase{\getDoctype{}} is my own work and I have documented all sources and material used.

\vspace{15mm}
\noindent
\getSubmissionLocation{}, \getSubmissionDate{} \hspace{50mm} \getAuthor{}

\cleardoublepage{}

%%%%%%%%%%%%%%%%%%%%%%%%%%%%%%%%%%%%%%%%%%%%%%%%%%%%%%%%%%%%%%
%                                                            %
%                                                            %
%                     ACKNOWLEDGMENTS                        %
%                                                            %
%                                                            %
%%%%%%%%%%%%%%%%%%%%%%%%%%%%%%%%%%%%%%%%%%%%%%%%%%%%%%%%%%%%%%

%\input{pages/acknowledgments}
\addcontentsline{toc}{chapter}{Acknowledgments}
\thispagestyle{empty}

\vspace*{20mm}

\begin{center}
{\usekomafont{section} Acknowledgments}
\end{center}

\vspace{10mm}

I owe an enormous debt of gratitude to my family which always suported me throughout my studies.
Thank you. I also would like to thank Prof. Tobias Nipkow for introducing me to the world of formal
verification through Isabelle and for supervising both my Bachelor's and my Master's thesis. It was
a pleasure to learn from and to work with you.

\cleardoublepage{}

%%%%%%%%%%%%%%%%%%%%%%%%%%%%%%%%%%%%%%%%%%%%%%%%%%%%%%%%%%%%%%
%                                                            %
%                                                            %
%                         ABSTRACT                           %
%                                                            %
%                                                            %
%%%%%%%%%%%%%%%%%%%%%%%%%%%%%%%%%%%%%%%%%%%%%%%%%%%%%%%%%%%%%%

\chapter{\abstractname}

In 1968 Earley \cite{Earley:1970} introduced his parsing algorithm capable of parsing all context-free grammars in cubic
space and time. This thesis presents the first formalization of an executable Earley parser in the interactive
theorem prover Isabelle/HOL \cite{Nipkow:2002}. We base our development on Jones' \cite{Jones:1972}
extensive paper proof of Earley's recognizer and the formalization of context-free grammars
and derivations of Obua \cite{Obua:2017}. We implement and prove correct a functional recognizer modeling Earley's
original imperative implementation and extend it with the necessary data structures to enable the construction
of parse trees following the work of Scott \cite{Scott:2008}. We then develop a functional algorithm that
builds a single parse tree and prove its correctness. Finally, we generalize this approach to an algorithm
for a complete parse forest.

\microtypesetup{protrusion=false}
\tableofcontents{}
\microtypesetup{protrusion=true}

\mainmatter{}

\endinput
%:%file=Limit.tex%:%

\endinput
%:%file=CFG.tex%:%

\endinput
%:%file=Derivations.tex%:%

\endinput
%:%file=Earley_Set.tex%:%

\endinput
%:%file=Earley_List.tex%:%

\endinput
%:%file=Examples.tex%:%

%
\begin{isabellebody}%
\setisabellecontext{{\isadigit{0}}{\isadigit{1}}{\isacharunderscore}{\kern0pt}Introduction}%
%
\isadelimtheory
%
\endisadelimtheory
%
\isatagtheory
%
\endisatagtheory
{\isafoldtheory}%
%
\isadelimtheory
\isanewline
%
\endisadelimtheory
\isacommand{notation}\isamarkupfalse%
\ {\isacharparenleft}{\kern0pt}latex\ \isakeyword{output}{\isacharparenright}{\kern0pt}\isanewline
\ \ nth\ \ {\isacharparenleft}{\kern0pt}{\isachardoublequoteopen}{\isacharunderscore}{\kern0pt}\ensuremath{_{[\mathit{{\isacharunderscore}{\kern0pt}}]}}{\isachardoublequoteclose}\ {\isacharbrackleft}{\kern0pt}{\isadigit{1}}{\isadigit{0}}{\isadigit{0}}{\isadigit{0}}{\isacharcomma}{\kern0pt}{\isadigit{0}}{\isacharbrackright}{\kern0pt}\ {\isadigit{1}}{\isadigit{0}}{\isadigit{0}}{\isadigit{0}}{\isacharparenright}{\kern0pt}%
\isadelimdocument
%
\endisadelimdocument
%
\isatagdocument
%
\isamarkupchapter{Questions:%
}
\isamarkuptrue%
%
\endisatagdocument
{\isafolddocument}%
%
\isadelimdocument
%
\endisadelimdocument
%
\begin{isamarkuptext}%
\begin{itemize}
  \item How long? Min/Max?
  \item How to define Earley inductively?
  \item Nicer notation for all the different kinds of derivations?
  \item Fix bibliography and Isabelle style???
\end{itemize}%
\end{isamarkuptext}\isamarkuptrue%
%
\isadelimdocument
%
\endisadelimdocument
%
\isatagdocument
%
\isamarkupchapter{Introduction%
}
\isamarkuptrue%
%
\isamarkupsection{Motivation%
}
\isamarkuptrue%
%
\endisatagdocument
{\isafolddocument}%
%
\isadelimdocument
%
\endisadelimdocument
%
\begin{isamarkuptext}%
some introduction about parsing, formal development of correct algorithms: an example based on
earley's recogniser, the benefits of formal methods, LocalLexing and the Bachelor thesis.%
\end{isamarkuptext}\isamarkuptrue%
%
\isadelimdocument
%
\endisadelimdocument
%
\isatagdocument
%
\isamarkupsection{Related Work%
}
\isamarkuptrue%
%
\endisatagdocument
{\isafolddocument}%
%
\isadelimdocument
%
\endisadelimdocument
%
\begin{isamarkuptext}%
Tomita \cite{Tomita:1987} presents an generalized LR parsing algorithm for augmented
context-free grammars that can handle arbitrary context-free grammars.

Izmaylova \textit{et al} \cite{Izmaylova:2016} develop a general parser 
combinator library based on memoized Continuation-Passing Style (CPS) recognizers that supports all
context-free grammars and constructs a Shared Packed Parse Forest (SPPF) in worst case cubic time and space.%
\end{isamarkuptext}\isamarkuptrue%
%
\begin{isamarkuptext}%
Obua \textit{et al} \cite{Obua:2017} introduce local lexing, a novel parsing concept which interleaves
lexing and parsing whilst allowing lexing to be dependent on the parsing process. They base their
development on Earley's algorithm and have verified the correctness with respect to its local lexing
semantics in the theorem prover Isabelle/HOL. The background theory of this Master's thesis is based
upon the local lexing entry \cite{LocalLexing-AFP} in the Archive of Formal Proofs.

Lasser \textit{et al} \cite{Lasser:2019} verify an LL(1) parser generator using the Coq proof assistant.

Barthwal \textit{et al} \cite{Barthwal:2009} formalize background theory
about context-free languages and grammars, and subsequently verify an SLR automaton and parser produced
by a parser generator.

Blaudeau \textit{et al} \cite{Blaudeau:2020} formalize the metatheory on Parsing expression grammars (PEGs) and
build a verified parser interpreter based on higher-order parsing combinators for expression grammars
using the PVS specification language and verification system. Koprowski \textit{et al} \cite{Koprowski:2011}
present TRX: a parser interpreter formally developed in Coq which also parses expression grammars.

Jourdan \textit{et al} \cite{Jourdan:2012} present a validator which checks if a context-free grammar
and an LR(1) parser agree, producing correctness guarantees required by verified compilers.

Lasser \textit{et al} \cite{Lasser:2021} present the verified parser CoStar based on the ALL(*) algorithm.
They proof soundness and completeness for all non-left-recursive grammars using the Coq proof assistant.%
\end{isamarkuptext}\isamarkuptrue%
%
\isadelimdocument
%
\endisadelimdocument
%
\isatagdocument
%
\isamarkupsection{Structure%
}
\isamarkuptrue%
%
\isamarkupsection{Contributions%
}
\isamarkuptrue%
%
\endisatagdocument
{\isafolddocument}%
%
\isadelimdocument
%
\endisadelimdocument
%
\begin{isamarkuptext}%
SNIPPET:

Context-free grammars have been used extensively for describing the syntax of programming languages
and natural languages. Parsing algorithms for context-free grammars consequently play a large role in
the implementation of compilers and interpreters for programming languages and of programs which understand
or translate natural languages. Numerous parsing algorithms have been developed. Some are general,
in the sense that they can handle all context-free grammars, while others can handle only subclasses of
grammars. The latter, restricted algorithms tend to be much more efficient The algorithm described here
seems to be the most efficient of the general algorithms, and also it can handle a larger class of grammars
in linear time than most of the restricted algorithms.%
\end{isamarkuptext}\isamarkuptrue%
%
\begin{isamarkuptext}%
SNIPPET:

The Computer Science community has been able to automatically generate parsers for a very wide class of context
free languages. However, many parsers are still written manually, either using tool support or even completely
by hand. This is partly because in some application areas such as natural language processing and bioinformatics
we don not have the luxury of designing the language so that it is amendable to know parsing techniques, but also
it is clear that left to themselves computer language designers do not naturally write LR(1) grammars. A grammar
not only defines the syntax of a language, it is also the starting point for the definition of the semantics,
and the grammar which facilitates semantics definition is not usually the one which is LR(1). Given this difficulty
in constructing natural LR(1) grammars that support desired semantics, the general parsing techniques, such as
the CYK Younger \cite{Younger:1967}, Earley \cite{Earley:1970} and GLR Tomita \cite{Tomita:1985} algorithms, developed
for natural language processing are also of interest to the wider computer science community. When using grammars as
the starting point for semantics definition, we distinguish between recognizers which simply determine whether or not
a given string is in the language defined by a given grammar, and parserwhich also return some form of derivation
of the string, if one exists. In their basic form the CYK and Earley algorithms are recognizers while GLR-style
algorithms are designed with derivation tree construction, and hence parsing, in mind.

There is no known liner time parsing or recognition algorithm that can be used with all context free grammars.
In their recognizer forms the CYK algorithm is worst case cubic on grammars in Chomsky normal form and Earley's
algorithm is worst case cubic on general context free grammers and worst case n2 on non-ambibuous grammars.
General recognizers must, by definition, be applicable to ambiguous grammars. Tomita's GLR algorithm is of unbounded
polynomial order in the worst case. Expanding general recognizers to parser raises several problems, not least
because there can be exponentially many or even infinitely many derivations for a given input string. A cubic
recognizer which was modified to simply return all derivations could become an unbounded parser.
Of course, it can be argued that ambiguous grammars reflect ambiguous semantics and thus should not be used in
practice. This would be far too extreme a position to take. For example, it is well known that the if-else
statement in hthe AnSI-standard grammar for C is ambiguous, but a longest match resolution results in a linear
time parser that attach the else to the most recent if, as specified by the ANSI-C semantics. The ambiguous
ANSI-C grammar is certainly practical for parser implementation. However, in general ambiguity is not so easily handled,
and it is well known that grammar ambiguity is in fact undecidable Hopcroft \textit{et al} \cite{Hopcroft:2006}, thus
we cannot expect a parser generator simply to check for ambiguity inthe grammar and report the problem back to the user.
Another possiblity is to avoid the issue by just returning one derivation. However, if only one derivation is returned
then this creates problems for a user who wants all derivations and, even in the case where only one derivation is
required, there is the issue of ensuring that it is the required derivationthat is returned. A truely general parser
will reutrn all possible derivations in some form. Perhaps the most well known representation is the shared packed
parse foreset SPPF described and used by Tomita \cite{Tomita:1985}. Tomita's description of the representation does ont allow
for the infinitely many derivations which arise from grammars which contain cycles, the source adapt the SPPF representation
to allow these. Johnson \cite{Johnson:1991} has shown that Tomita-style SPPFs are worst case unbounded polynomial size. Thus
using such structures will alo turn any cubic recognition technique into a worst case unbounded polynomial parsing technique.
Leaving aside the potential increase in complexity when turning a recogniser into a parser, it is clear that this proccess is often difficult to carry
out correctly. Earley gave an algorithm for constructing derivations of a string accepted by his recognizer,
but this was subsequently shown by Tomita \cite{Tomita:1985} to return spurious derivations in certain cases.
Tomita's original version of his algorithm failed to terminate on grammars with hidden left recursio and, as remarked above
, had no mechanism for contructing complete SPPFs for grammers with cycles.%
\end{isamarkuptext}\isamarkuptrue%
%
\isadelimtheory
%
\endisadelimtheory
%
\isatagtheory
%
\endisatagtheory
{\isafoldtheory}%
%
\isadelimtheory
%
\endisadelimtheory
%
\end{isabellebody}%
\endinput
%:%file=01_Introduction.tex%:%
%:%15=7%:%
%:%18=8%:%
%:%19=8%:%
%:%20=9%:%
%:%27=11%:%
%:%39=14%:%
%:%40=15%:%
%:%41=16%:%
%:%42=17%:%
%:%43=18%:%
%:%44=19%:%
%:%53=22%:%
%:%57=24%:%
%:%69=26%:%
%:%70=27%:%
%:%79=29%:%
%:%91=31%:%
%:%92=32%:%
%:%93=33%:%
%:%94=34%:%
%:%95=35%:%
%:%96=36%:%
%:%100=40%:%
%:%101=41%:%
%:%102=42%:%
%:%103=43%:%
%:%104=44%:%
%:%105=45%:%
%:%106=46%:%
%:%107=47%:%
%:%108=48%:%
%:%109=49%:%
%:%110=50%:%
%:%111=51%:%
%:%112=52%:%
%:%113=53%:%
%:%114=54%:%
%:%115=55%:%
%:%116=56%:%
%:%117=57%:%
%:%118=58%:%
%:%119=59%:%
%:%120=60%:%
%:%121=61%:%
%:%130=64%:%
%:%134=66%:%
%:%146=69%:%
%:%147=70%:%
%:%148=71%:%
%:%149=72%:%
%:%150=73%:%
%:%151=74%:%
%:%152=75%:%
%:%153=76%:%
%:%154=77%:%
%:%155=78%:%
%:%159=82%:%
%:%160=83%:%
%:%161=84%:%
%:%162=85%:%
%:%163=86%:%
%:%164=87%:%
%:%165=88%:%
%:%166=89%:%
%:%167=90%:%
%:%168=91%:%
%:%169=92%:%
%:%170=93%:%
%:%171=94%:%
%:%172=95%:%
%:%173=96%:%
%:%174=97%:%
%:%175=98%:%
%:%176=99%:%
%:%177=100%:%
%:%178=101%:%
%:%179=102%:%
%:%180=103%:%
%:%181=104%:%
%:%182=105%:%
%:%183=106%:%
%:%184=107%:%
%:%185=108%:%
%:%186=109%:%
%:%187=110%:%
%:%188=111%:%
%:%189=112%:%
%:%190=113%:%
%:%191=114%:%
%:%192=115%:%
%:%193=116%:%
%:%194=117%:%
%:%195=118%:%
%:%196=119%:%
%:%197=120%:%
%:%198=121%:%
%:%199=122%:%
%:%200=123%:%
%:%201=124%:%
%:%202=125%:%

%
\begin{isabellebody}%
\setisabellecontext{{\isadigit{0}}{\isadigit{2}}{\isacharunderscore}{\kern0pt}Earleys{\isacharunderscore}{\kern0pt}Algorithm}%
%
\isadelimtheory
%
\endisadelimtheory
%
\isatagtheory
%
\endisatagtheory
{\isafoldtheory}%
%
\isadelimtheory
%
\endisadelimtheory
%
\isadelimdocument
%
\endisadelimdocument
%
\isatagdocument
%
\isamarkupchapter{Earley's Algorithm%
}
\isamarkuptrue%
%
\endisatagdocument
{\isafolddocument}%
%
\isadelimdocument
%
\endisadelimdocument
%
\begin{isamarkuptext}%
We present a slightly simplified version of Earley's original recognizer algorithm \cite{Earley:1970},
omitting Earley's proposed look-ahead since it's primary purpose is to increase the efficiency of the
resulting recognizer. Throughout this thesis we are working with a running example. The considered grammar is a tiny excerpt of a toy
arithmetic expression grammar: \isa{{\isasymG}} $::= S \rightarrow \, x \, \vert \, S \rightarrow \, S + S$ and
the input is \isa{{\isasymomega}} $= x + x + x$.

Intuitively, Earley's recognizer works in principle like a top-down parser carrying along all possible
parses simultaneously in an efficient manner.
In detail, the algorithm works as follows: it scans the input \isa{{\isasymomega}} $=a_0,\dots,a_n$, constructing
$n+1$ Earley bins $B_i$ which are sets of Earley items. An inital bin $B_0$ and one bin $B_{i+1}$ for
each symbol $a_i$ of the input.
In general, an Earley item $A \rightarrow \, \alpha \bullet \beta, i, j$ consists of four parts: a production rule of the grammar which we are currently
scanning, a bullet signalling how much of the production's right-hand side we have recognized so far,
an origin $i$ describing the position of \isa{{\isasymomega}} where we started scanning, and an end $j$ indicating
the position  of \isa{{\isasymomega}} we are currently considering next for the remaining right-hand side of the production rule.
Note that there will be only one set of earley items or only one bin $B$ and we say an item is conceptually part of bin $B_j$ if it's end is the index $j$.
Table \ref{tab:earley_bins} lists the items for our example grammar. Bin $B_4$ contains for example the item $S \rightarrow \, S + \bullet S, 2, 4$.
Or, we are scanning the rule $S \rightarrow \, S + S$, have recognized the substring from $2$ to $4$ (the first index being
inclusive the second one exclusive) of \isa{{\isasymomega}} by $\alpha = S +$, and are trying to scan $\beta = S$ from position 4 in \omega. 

The algorithm initializes $B$ by applying the \textit{Init} operation. It then proceeds to execute
the \textit{Scan}, \textit{Predict} and \textit{Complete} operations listed in figure \ref{fig:inference_rules}
until there are no more new items being generated and added to $B$. Next we describe these four operations
in detail:

\begin{enumerate}
  \item The \textit{Init} operation adds items
    $S \rightarrow \, \bullet\alpha, 0, 0$ for each production rule containing the start symbol $S$ on its left-hand side.

    For our example \textit{Init} adds the items $S \rightarrow \, \bullet x, 0, 0$ and $S \rightarrow \, \bullet S + S, 0 , 0$.
  \item The \textit{Scan} operation applies if there is a terminal to the right side of the bullet, or items of the form $A \rightarrow \, \alpha \bullet a \beta, i, j$,
    and the $j$-th symbol of \isa{{\isasymomega}} matches the terminal symbol following the bullet. We add one new item $A \rightarrow \, \alpha a \bullet \beta, i, j+1$
    to $B$ moving the bullet over the scanned terminal symbol.

    Considering our example, bin $B_3$ contains
    the item $S \rightarrow \, S \bullet + S, 2, 3$, the third element of \isa{{\isasymomega}} is the terminal $+$, so we add the
    item $S \rightarrow \, S + \bullet S, 2, 4$ to the conceptual bin $B_4$.
  \item The \textit{Predict} operation is applicable to an item when there is a non-terminal to the right of
    the bullet or items of the form $A \rightarrow \, \alpha \bullet B \beta, i, j$. It adds one new item $B \rightarrow \, \bullet \gamma, j, j$
    to the bin for each alternate $B \rightarrow \, \gamma$ of that non-terminal.

    E.g. for the item  $S \rightarrow \, S + \bullet S, 0, 2$ in $B_2$
    we add the two items $S \rightarrow \, \bullet x, 2, 2$ and $S \rightarrow \, \bullet S + S, 2, 2$ corresponding
    to the two alternates of $S$. The bullet is set to the beginning of the right-hand side of the production
    rule, the origin and end are set to $j = 2$ to indicate that we are starting to scan in the current bin and
    have not scanned anything so far.
  \item The \textit{Complete} operation applies if we process an item with the bullet at the end of the
    right side of its production rule. For an item $B \rightarrow \, \gamma \bullet, j, k$ we have successfully scanned the substring
    \omega[j..k) and are now going back to the origin bin $B_j$ where we predicted this non-terminal. There we look for any item of the form
    $A \rightarrow \, \alpha \bullet B \beta, i, j$ containing a bullet in front of the non-terminal we completed, or the reason we
    predicted it on the first place. Since we scanned the predicted non-terminal successfully, we are allowed to
    move over the bullet, resulting in one new item $A \rightarrow \, \alpha B \bullet \beta, i, k$. Note in particular
    the origin and end indices.

    Looking back at our example, we can add the item $S \rightarrow \, S + S \bullet, 0, 5$
    for two different reasons corresponding conceptually to the two different ways we can derive \omega.
    When processing $S \rightarrow \, x \bullet, 4, 5$ we find $S \rightarrow \, S + \bullet S, 0, 4$ in the origin
    bin $B_4$ which conceptually corresponds to recognizing $(x + x) + x$. We \glq add \grq the same item again
    while applying the \textit{Complete} operation to $S \rightarrow \, S + S \bullet, 2, 5$ and $S \rightarrow \, S + \bullet S, 0, 2$
    which corresponds to recognizing the input as $x + (x + x)$.
\end{enumerate}

If the algorithm encounters an item of the form $S \rightarrow \, \alpha, 0, \isa{{\isacharbar}{\kern0pt}{\isasymomega}{\isacharbar}{\kern0pt}\ {\isacharplus}{\kern0pt}\ {\isadigit{1}}}$,
it returns \textit{true}, otherwise it returns \textit{false}. For the tiny arithmetic expression grammar
we generate the item $S \rightarrow \, S + S \bullet, 0, 5$ and return the correct answer \textit{true},
since there exist derivations for $\omega = x + x + x$, e.g.
$S \Rightarrow S + S \Rightarrow x + S \Rightarrow x + S + S \xRightarrow{\ast} x + x + x$ or
$S \Rightarrow S + S \Rightarrow S + x \Rightarrow S + S + x \xRightarrow{\ast} x + x + x$.

To proof the correctness of Earley's recognizer algorithm we need to show the following theorem:

$$S \rightarrow \, \alpha \bullet, 0, \isa{{\isacharbar}{\kern0pt}{\isasymomega}{\isacharbar}{\kern0pt}\ {\isacharplus}{\kern0pt}\ {\isadigit{1}}} \in B \,\,\, \textrm{iff} \,\,\, S \, \xRightarrow{\ast} \, \isa{{\isasymomega}}$$

It follows from the following three lemmas:

\begin{enumerate}
  \item Termination: there only exist a finite number of Earley items
  \item Soundness: for every generated item there exists an according derivation: \\
     $A \rightarrow \, \alpha \bullet \beta, i, j \in B \,\,\, \textrm{implies} \,\,\, A \, \xRightarrow{\ast} \, \isa{{\isasymomega}}[i..j)$
  \item Completeness: for every derivation we generate an according item: \\
     $A \, \xRightarrow{\ast} \, \isa{{\isasymomega}}[i..j) \,\,\, \textrm{implies} \,\,\, A \rightarrow \, \alpha \bullet \beta, i, j \in B$
\end{enumerate}%
\end{isamarkuptext}\isamarkuptrue%
%
\begin{isamarkuptext}%
\begin{figure}[htpb]
    \centering

    \begin{mathpar}
      \inferrule [Init]
      {\\}
      {$S \rightarrow \, \bullet\alpha, 0, 0$}
  
      \inferrule [Scan]
      {$A \rightarrow \, \alpha \bullet a \beta, i, j$ \\ $\isa{{\isasymomega}}[j] = a$}
      {$A \rightarrow \, \alpha a \bullet \beta, i, j+1$}
  
      \inferrule [Predict]
      {$A \rightarrow \, \alpha \bullet B \beta, i, j$ \\ $B \rightarrow \, \gamma \, \in \, \isa{{\isasymG}}$}
      {$B \rightarrow \, \bullet \gamma, j, j$}
  
      \inferrule [Complete]
      {$A \rightarrow \, \alpha \bullet B \beta, i, j$ \\ $B \rightarrow \, \gamma \bullet, j, k$}
      {$A \rightarrow \, \alpha B \bullet \beta, i, k$}
    \end{mathpar}
    \caption[Earley inference rules]{Earley inference rules}\label{fig:earley-inference-rules}
    \label{fig:inference_rules}
  \end{figure}%
\end{isamarkuptext}\isamarkuptrue%
%
\begin{isamarkuptext}%
\begin{table}[htpb] 
    \caption[Earley items running example]{Earley items for the grammar \isa{{\isasymG}}: $S \rightarrow \, x$, $S \rightarrow \, S + S$}\label{tab:earley-items}
    \centering
    \begin{tabular}{| l | l | l |}
        $B_0$                                   & $B_1$                                    & $B_2$                                \\
      \midrule
        $S \rightarrow \, \bullet x, 0, 0$      & $S \rightarrow \, x \bullet, 0, 1$     & $S \rightarrow \, S + \bullet S, 0, 2$ \\
        $S \rightarrow \, \bullet S + S, 0 , 0$ & $S \rightarrow \, S \bullet + S, 0, 1$ & $S \rightarrow \, \bullet x, 2, 2$     \\
                                                &                                        & $S \rightarrow \, \bullet S + S, 2, 2$ \\

      \midrule

        $B_3$                                  & $B_4$                                    & $B_5$                                \\
      \midrule
        $S \rightarrow \, x \bullet, 2, 3$     & $S \rightarrow \, S + \bullet S, 2, 4$ & $S \rightarrow \, x \bullet, 4, 5$     \\
        $S \rightarrow \, S + S \bullet, 0, 3$ & $S \rightarrow \, S + \bullet S, 0, 4$ & $S \rightarrow \, S + S \bullet, 2, 5$ \\
        $S \rightarrow \, S \bullet + S, 2, 3$ & $S \rightarrow \, \bullet x, 4, 4$     & $S \rightarrow \, S + S \bullet, 0, 5$ \\
        $S \rightarrow \, S \bullet + S, 0, 3$ & $S \rightarrow \, \bullet S + S, 4, 4$ & $S \rightarrow \, S \bullet + S, 4, 5$ \\
                                               &                                        & $S \rightarrow \, S \bullet + S, 2, 5$ \\
                                               &                                        & $S \rightarrow \, S \bullet + S, 0, 5$ \\
    \end{tabular}
    \label{tab:earley_bins}
  \end{table}%
\end{isamarkuptext}\isamarkuptrue%
%
\isadelimtheory
%
\endisadelimtheory
%
\isatagtheory
%
\endisatagtheory
{\isafoldtheory}%
%
\isadelimtheory
%
\endisadelimtheory
%
\end{isabellebody}%
\endinput
%:%file=02_Earleys_Algorithm.tex%:%
%:%24=8%:%
%:%36=11%:%
%:%37=12%:%
%:%38=13%:%
%:%39=14%:%
%:%40=15%:%
%:%41=16%:%
%:%42=17%:%
%:%43=18%:%
%:%44=19%:%
%:%45=20%:%
%:%46=21%:%
%:%47=22%:%
%:%48=23%:%
%:%49=24%:%
%:%50=25%:%
%:%51=26%:%
%:%52=27%:%
%:%53=28%:%
%:%54=29%:%
%:%55=30%:%
%:%56=31%:%
%:%57=32%:%
%:%58=33%:%
%:%59=34%:%
%:%60=35%:%
%:%61=36%:%
%:%62=37%:%
%:%63=38%:%
%:%64=39%:%
%:%65=40%:%
%:%66=41%:%
%:%67=42%:%
%:%68=43%:%
%:%69=44%:%
%:%70=45%:%
%:%71=46%:%
%:%72=47%:%
%:%73=48%:%
%:%74=49%:%
%:%75=50%:%
%:%76=51%:%
%:%77=52%:%
%:%78=53%:%
%:%79=54%:%
%:%80=55%:%
%:%81=56%:%
%:%82=57%:%
%:%83=58%:%
%:%84=59%:%
%:%85=60%:%
%:%86=61%:%
%:%87=62%:%
%:%88=63%:%
%:%89=64%:%
%:%90=65%:%
%:%91=66%:%
%:%92=67%:%
%:%93=68%:%
%:%94=69%:%
%:%95=70%:%
%:%96=71%:%
%:%97=72%:%
%:%98=73%:%
%:%99=74%:%
%:%100=75%:%
%:%101=76%:%
%:%102=77%:%
%:%103=78%:%
%:%104=79%:%
%:%105=80%:%
%:%106=81%:%
%:%107=82%:%
%:%108=83%:%
%:%109=84%:%
%:%110=85%:%
%:%111=86%:%
%:%112=87%:%
%:%113=88%:%
%:%114=89%:%
%:%115=90%:%
%:%116=91%:%
%:%117=92%:%
%:%121=96%:%
%:%122=97%:%
%:%123=98%:%
%:%124=99%:%
%:%125=100%:%
%:%126=101%:%
%:%127=102%:%
%:%128=103%:%
%:%129=104%:%
%:%130=105%:%
%:%131=106%:%
%:%132=107%:%
%:%133=108%:%
%:%134=109%:%
%:%135=110%:%
%:%136=111%:%
%:%137=112%:%
%:%138=113%:%
%:%139=114%:%
%:%140=115%:%
%:%141=116%:%
%:%142=117%:%
%:%143=118%:%
%:%147=122%:%
%:%148=123%:%
%:%149=124%:%
%:%150=125%:%
%:%151=126%:%
%:%152=127%:%
%:%153=128%:%
%:%154=129%:%
%:%155=130%:%
%:%156=131%:%
%:%157=132%:%
%:%158=133%:%
%:%159=134%:%
%:%160=135%:%
%:%161=136%:%
%:%162=137%:%
%:%163=138%:%
%:%164=139%:%
%:%165=140%:%
%:%166=141%:%
%:%167=142%:%
%:%168=143%:%
%:%169=144%:%

%
\begin{isabellebody}%
\setisabellecontext{{\isadigit{0}}{\isadigit{3}}{\isacharunderscore}{\kern0pt}Fixpoint{\isacharunderscore}{\kern0pt}Earley{\isacharunderscore}{\kern0pt}Recogniser}%
%
\isadelimtheory
%
\endisadelimtheory
%
\isatagtheory
%
\endisatagtheory
{\isafoldtheory}%
%
\isadelimtheory
%
\endisadelimtheory
%
\isadelimdocument
%
\endisadelimdocument
%
\isatagdocument
%
\isamarkupchapter{Fixpoint Recogniser%
}
\isamarkuptrue%
%
\endisatagdocument
{\isafolddocument}%
%
\isadelimdocument
%
\endisadelimdocument
%
\isadelimtheory
%
\endisadelimtheory
%
\isatagtheory
%
\endisatagtheory
{\isafoldtheory}%
%
\isadelimtheory
%
\endisadelimtheory
%
\end{isabellebody}%
\endinput
%:%file=03_Fixpoint_Earley_Recogniser.tex%:%
%:%24=9%:%

%
\begin{isabellebody}%
\setisabellecontext{{\isadigit{0}}{\isadigit{4}}{\isacharunderscore}{\kern0pt}Earley{\isacharunderscore}{\kern0pt}Recogniser}%
%
\isadelimtheory
%
\endisadelimtheory
%
\isatagtheory
%
\endisatagtheory
{\isafoldtheory}%
%
\isadelimtheory
%
\endisadelimtheory
%
\isadelimdocument
%
\endisadelimdocument
%
\isatagdocument
%
\isamarkupchapter{Earley Recogniser%
}
\isamarkuptrue%
%
\endisatagdocument
{\isafolddocument}%
%
\isadelimdocument
%
\endisadelimdocument
%
\isadelimtheory
%
\endisadelimtheory
%
\isatagtheory
%
\endisatagtheory
{\isafoldtheory}%
%
\isadelimtheory
%
\endisadelimtheory
%
\end{isabellebody}%
\endinput
%:%file=04_Earley_Recogniser.tex%:%
%:%24=9%:%

%
\begin{isabellebody}%
\setisabellecontext{{\isadigit{0}}{\isadigit{5}}{\isacharunderscore}{\kern0pt}Earley{\isacharunderscore}{\kern0pt}Parser}%
%
\isadelimtheory
%
\endisadelimtheory
%
\isatagtheory
%
\endisatagtheory
{\isafoldtheory}%
%
\isadelimtheory
%
\endisadelimtheory
%
\isadelimdocument
%
\endisadelimdocument
%
\isatagdocument
%
\isamarkupchapter{Earley Parser Implementation \label{chap:05}%
}
\isamarkuptrue%
%
\isamarkupsection{A Single Parse Tree%
}
\isamarkuptrue%
%
\endisatagdocument
{\isafolddocument}%
%
\isadelimdocument
%
\endisadelimdocument
\isacommand{datatype}\isamarkupfalse%
\ {\isacharprime}{\kern0pt}a\ tree\ {\isacharequal}{\kern0pt}\isanewline
\ \ Leaf\ {\isacharprime}{\kern0pt}a\isanewline
\ \ {\isacharbar}{\kern0pt}\ Branch\ {\isacharprime}{\kern0pt}a\ {\isachardoublequoteopen}{\isacharprime}{\kern0pt}a\ tree\ list{\isachardoublequoteclose}\isanewline
\isanewline
\isacommand{fun}\isamarkupfalse%
\ yield{\isacharunderscore}{\kern0pt}tree\ {\isacharcolon}{\kern0pt}{\isacharcolon}{\kern0pt}\ {\isachardoublequoteopen}{\isacharprime}{\kern0pt}a\ tree\ {\isasymRightarrow}\ {\isacharprime}{\kern0pt}a\ sentential{\isachardoublequoteclose}\ \isakeyword{where}\isanewline
\ \ {\isachardoublequoteopen}yield{\isacharunderscore}{\kern0pt}tree\ {\isacharparenleft}{\kern0pt}Leaf\ a{\isacharparenright}{\kern0pt}\ {\isacharequal}{\kern0pt}\ {\isacharbrackleft}{\kern0pt}a{\isacharbrackright}{\kern0pt}{\isachardoublequoteclose}\isanewline
{\isacharbar}{\kern0pt}\ {\isachardoublequoteopen}yield{\isacharunderscore}{\kern0pt}tree\ {\isacharparenleft}{\kern0pt}Branch\ {\isacharunderscore}{\kern0pt}\ ts{\isacharparenright}{\kern0pt}\ {\isacharequal}{\kern0pt}\ concat\ {\isacharparenleft}{\kern0pt}map\ yield{\isacharunderscore}{\kern0pt}tree\ ts{\isacharparenright}{\kern0pt}{\isachardoublequoteclose}\isanewline
\isanewline
\isacommand{fun}\isamarkupfalse%
\ root{\isacharunderscore}{\kern0pt}tree\ {\isacharcolon}{\kern0pt}{\isacharcolon}{\kern0pt}\ {\isachardoublequoteopen}{\isacharprime}{\kern0pt}a\ tree\ {\isasymRightarrow}\ {\isacharprime}{\kern0pt}a{\isachardoublequoteclose}\ \isakeyword{where}\isanewline
\ \ {\isachardoublequoteopen}root{\isacharunderscore}{\kern0pt}tree\ {\isacharparenleft}{\kern0pt}Leaf\ a{\isacharparenright}{\kern0pt}\ {\isacharequal}{\kern0pt}\ a{\isachardoublequoteclose}\isanewline
{\isacharbar}{\kern0pt}\ {\isachardoublequoteopen}root{\isacharunderscore}{\kern0pt}tree\ {\isacharparenleft}{\kern0pt}Branch\ N\ {\isacharunderscore}{\kern0pt}{\isacharparenright}{\kern0pt}\ {\isacharequal}{\kern0pt}\ N{\isachardoublequoteclose}\isanewline
\isanewline
\isacommand{fun}\isamarkupfalse%
\ wf{\isacharunderscore}{\kern0pt}rule{\isacharunderscore}{\kern0pt}tree\ {\isacharcolon}{\kern0pt}{\isacharcolon}{\kern0pt}\ {\isachardoublequoteopen}{\isacharprime}{\kern0pt}a\ cfg\ {\isasymRightarrow}\ {\isacharprime}{\kern0pt}a\ tree\ {\isasymRightarrow}\ bool{\isachardoublequoteclose}\ \isakeyword{where}\isanewline
\ \ {\isachardoublequoteopen}wf{\isacharunderscore}{\kern0pt}rule{\isacharunderscore}{\kern0pt}tree\ {\isacharunderscore}{\kern0pt}\ {\isacharparenleft}{\kern0pt}Leaf\ a{\isacharparenright}{\kern0pt}\ {\isasymlongleftrightarrow}\ True{\isachardoublequoteclose}\isanewline
{\isacharbar}{\kern0pt}\ {\isachardoublequoteopen}wf{\isacharunderscore}{\kern0pt}rule{\isacharunderscore}{\kern0pt}tree\ {\isasymG}\ {\isacharparenleft}{\kern0pt}Branch\ N\ ts{\isacharparenright}{\kern0pt}\ {\isasymlongleftrightarrow}\ {\isacharparenleft}{\kern0pt}\isanewline
\ \ \ \ {\isacharparenleft}{\kern0pt}{\isasymexists}r\ {\isasymin}\ set\ {\isacharparenleft}{\kern0pt}{\isasymRR}\ {\isasymG}{\isacharparenright}{\kern0pt}{\isachardot}{\kern0pt}\ N\ {\isacharequal}{\kern0pt}\ rule{\isacharunderscore}{\kern0pt}head\ r\ {\isasymand}\ map\ root{\isacharunderscore}{\kern0pt}tree\ ts\ {\isacharequal}{\kern0pt}\ rule{\isacharunderscore}{\kern0pt}body\ r{\isacharparenright}{\kern0pt}\ {\isasymand}\isanewline
\ \ \ \ {\isacharparenleft}{\kern0pt}{\isasymforall}t\ {\isasymin}\ set\ ts{\isachardot}{\kern0pt}\ wf{\isacharunderscore}{\kern0pt}rule{\isacharunderscore}{\kern0pt}tree\ {\isasymG}\ t{\isacharparenright}{\kern0pt}{\isacharparenright}{\kern0pt}{\isachardoublequoteclose}\isanewline
\isanewline
\isacommand{fun}\isamarkupfalse%
\ wf{\isacharunderscore}{\kern0pt}item{\isacharunderscore}{\kern0pt}tree\ {\isacharcolon}{\kern0pt}{\isacharcolon}{\kern0pt}\ {\isachardoublequoteopen}{\isacharprime}{\kern0pt}a\ cfg\ {\isasymRightarrow}\ {\isacharprime}{\kern0pt}a\ item\ {\isasymRightarrow}\ {\isacharprime}{\kern0pt}a\ tree\ {\isasymRightarrow}\ bool{\isachardoublequoteclose}\ \isakeyword{where}\isanewline
\ \ {\isachardoublequoteopen}wf{\isacharunderscore}{\kern0pt}item{\isacharunderscore}{\kern0pt}tree\ {\isasymG}\ {\isacharunderscore}{\kern0pt}\ {\isacharparenleft}{\kern0pt}Leaf\ a{\isacharparenright}{\kern0pt}\ {\isasymlongleftrightarrow}\ True{\isachardoublequoteclose}\isanewline
{\isacharbar}{\kern0pt}\ {\isachardoublequoteopen}wf{\isacharunderscore}{\kern0pt}item{\isacharunderscore}{\kern0pt}tree\ {\isasymG}\ x\ {\isacharparenleft}{\kern0pt}Branch\ N\ ts{\isacharparenright}{\kern0pt}\ {\isasymlongleftrightarrow}\ {\isacharparenleft}{\kern0pt}\isanewline
\ \ \ \ N\ {\isacharequal}{\kern0pt}\ item{\isacharunderscore}{\kern0pt}rule{\isacharunderscore}{\kern0pt}head\ x\ {\isasymand}\isanewline
\ \ \ \ map\ root{\isacharunderscore}{\kern0pt}tree\ ts\ {\isacharequal}{\kern0pt}\ take\ {\isacharparenleft}{\kern0pt}item{\isacharunderscore}{\kern0pt}bullet\ x{\isacharparenright}{\kern0pt}\ {\isacharparenleft}{\kern0pt}item{\isacharunderscore}{\kern0pt}rule{\isacharunderscore}{\kern0pt}body\ x{\isacharparenright}{\kern0pt}\ {\isasymand}\isanewline
\ \ \ \ {\isacharparenleft}{\kern0pt}{\isasymforall}t\ {\isasymin}\ set\ ts{\isachardot}{\kern0pt}\ wf{\isacharunderscore}{\kern0pt}rule{\isacharunderscore}{\kern0pt}tree\ {\isasymG}\ t{\isacharparenright}{\kern0pt}{\isacharparenright}{\kern0pt}{\isachardoublequoteclose}\isanewline
\isanewline
\isacommand{definition}\isamarkupfalse%
\ wf{\isacharunderscore}{\kern0pt}yield{\isacharunderscore}{\kern0pt}tree\ {\isacharcolon}{\kern0pt}{\isacharcolon}{\kern0pt}\ {\isachardoublequoteopen}{\isacharprime}{\kern0pt}a\ sentential\ {\isasymRightarrow}\ {\isacharprime}{\kern0pt}a\ item\ {\isasymRightarrow}\ {\isacharprime}{\kern0pt}a\ tree\ {\isasymRightarrow}\ bool{\isachardoublequoteclose}\ \isakeyword{where}\isanewline
\ \ {\isachardoublequoteopen}wf{\isacharunderscore}{\kern0pt}yield{\isacharunderscore}{\kern0pt}tree\ {\isasymomega}\ x\ t\ {\isasymequiv}\ yield{\isacharunderscore}{\kern0pt}tree\ t\ {\isacharequal}{\kern0pt}\ {\isasymomega}{\isacharbrackleft}{\kern0pt}item{\isacharunderscore}{\kern0pt}origin\ x{\isachardot}{\kern0pt}{\isachardot}{\kern0pt}item{\isacharunderscore}{\kern0pt}end\ x{\isasymrangle}{\isachardoublequoteclose}%
\isadelimdocument
%
\endisadelimdocument
%
\isatagdocument
%
\isamarkupsubsection{Pointer Lemmas%
}
\isamarkuptrue%
%
\endisatagdocument
{\isafolddocument}%
%
\isadelimdocument
%
\endisadelimdocument
\isacommand{definition}\isamarkupfalse%
\ predicts\ {\isacharcolon}{\kern0pt}{\isacharcolon}{\kern0pt}\ {\isachardoublequoteopen}{\isacharprime}{\kern0pt}a\ item\ {\isasymRightarrow}\ bool{\isachardoublequoteclose}\ \isakeyword{where}\isanewline
\ \ {\isachardoublequoteopen}predicts\ x\ {\isasymequiv}\ item{\isacharunderscore}{\kern0pt}origin\ x\ {\isacharequal}{\kern0pt}\ item{\isacharunderscore}{\kern0pt}end\ x\ {\isasymand}\ item{\isacharunderscore}{\kern0pt}bullet\ x\ {\isacharequal}{\kern0pt}\ {\isadigit{0}}{\isachardoublequoteclose}\isanewline
\isanewline
\isacommand{definition}\isamarkupfalse%
\ scans\ {\isacharcolon}{\kern0pt}{\isacharcolon}{\kern0pt}\ {\isachardoublequoteopen}{\isacharprime}{\kern0pt}a\ sentential\ {\isasymRightarrow}\ nat\ {\isasymRightarrow}\ {\isacharprime}{\kern0pt}a\ item\ {\isasymRightarrow}\ {\isacharprime}{\kern0pt}a\ item\ {\isasymRightarrow}\ bool{\isachardoublequoteclose}\ \isakeyword{where}\isanewline
\ \ {\isachardoublequoteopen}scans\ {\isasymomega}\ k\ x\ y\ {\isasymequiv}\ y\ {\isacharequal}{\kern0pt}\ inc{\isacharunderscore}{\kern0pt}item\ x\ k\ {\isasymand}\ {\isacharparenleft}{\kern0pt}{\isasymexists}a{\isachardot}{\kern0pt}\ next{\isacharunderscore}{\kern0pt}symbol\ x\ {\isacharequal}{\kern0pt}\ Some\ a\ {\isasymand}\ {\isasymomega}{\isacharbang}{\kern0pt}{\isacharparenleft}{\kern0pt}k{\isacharminus}{\kern0pt}{\isadigit{1}}{\isacharparenright}{\kern0pt}\ {\isacharequal}{\kern0pt}\ a{\isacharparenright}{\kern0pt}{\isachardoublequoteclose}\isanewline
\isanewline
\isacommand{definition}\isamarkupfalse%
\ completes\ {\isacharcolon}{\kern0pt}{\isacharcolon}{\kern0pt}\ {\isachardoublequoteopen}nat\ {\isasymRightarrow}\ {\isacharprime}{\kern0pt}a\ item\ {\isasymRightarrow}\ {\isacharprime}{\kern0pt}a\ item\ {\isasymRightarrow}\ {\isacharprime}{\kern0pt}a\ item\ {\isasymRightarrow}\ bool{\isachardoublequoteclose}\ \isakeyword{where}\isanewline
\ \ {\isachardoublequoteopen}completes\ k\ x\ y\ z\ {\isasymequiv}\ y\ {\isacharequal}{\kern0pt}\ inc{\isacharunderscore}{\kern0pt}item\ x\ k\ {\isasymand}\isanewline
\ \ \ \ is{\isacharunderscore}{\kern0pt}complete\ z\ {\isasymand}\isanewline
\ \ \ \ item{\isacharunderscore}{\kern0pt}origin\ z\ {\isacharequal}{\kern0pt}\ item{\isacharunderscore}{\kern0pt}end\ x\ {\isasymand}\isanewline
\ \ \ \ {\isacharparenleft}{\kern0pt}{\isasymexists}N{\isachardot}{\kern0pt}\ next{\isacharunderscore}{\kern0pt}symbol\ x\ {\isacharequal}{\kern0pt}\ Some\ N\ {\isasymand}\ N\ {\isacharequal}{\kern0pt}\ item{\isacharunderscore}{\kern0pt}rule{\isacharunderscore}{\kern0pt}head\ z{\isacharparenright}{\kern0pt}{\isachardoublequoteclose}\isanewline
\isanewline
\isacommand{definition}\isamarkupfalse%
\ sound{\isacharunderscore}{\kern0pt}null{\isacharunderscore}{\kern0pt}ptr\ {\isacharcolon}{\kern0pt}{\isacharcolon}{\kern0pt}\ {\isachardoublequoteopen}{\isacharprime}{\kern0pt}a\ entry\ {\isasymRightarrow}\ bool{\isachardoublequoteclose}\ \isakeyword{where}\isanewline
\ \ {\isachardoublequoteopen}sound{\isacharunderscore}{\kern0pt}null{\isacharunderscore}{\kern0pt}ptr\ e\ {\isasymequiv}\ pointer\ e\ {\isacharequal}{\kern0pt}\ Null\ {\isasymlongrightarrow}\ predicts\ {\isacharparenleft}{\kern0pt}item\ e{\isacharparenright}{\kern0pt}{\isachardoublequoteclose}\isanewline
\isanewline
\isacommand{definition}\isamarkupfalse%
\ sound{\isacharunderscore}{\kern0pt}pre{\isacharunderscore}{\kern0pt}ptr\ {\isacharcolon}{\kern0pt}{\isacharcolon}{\kern0pt}\ {\isachardoublequoteopen}{\isacharprime}{\kern0pt}a\ sentential\ {\isasymRightarrow}\ {\isacharprime}{\kern0pt}a\ bins\ {\isasymRightarrow}\ nat\ {\isasymRightarrow}\ {\isacharprime}{\kern0pt}a\ entry\ {\isasymRightarrow}\ bool{\isachardoublequoteclose}\ \isakeyword{where}\isanewline
\ \ {\isachardoublequoteopen}sound{\isacharunderscore}{\kern0pt}pre{\isacharunderscore}{\kern0pt}ptr\ {\isasymomega}\ bs\ k\ e\ {\isasymequiv}\ {\isasymforall}pre{\isachardot}{\kern0pt}\ pointer\ e\ {\isacharequal}{\kern0pt}\ Pre\ pre\ {\isasymlongrightarrow}\isanewline
\ \ \ \ k\ {\isachargreater}{\kern0pt}\ {\isadigit{0}}\ {\isasymand}\isanewline
\ \ \ \ pre\ {\isacharless}{\kern0pt}\ {\isacharbar}{\kern0pt}bs{\isacharbang}{\kern0pt}{\isacharparenleft}{\kern0pt}k{\isacharminus}{\kern0pt}{\isadigit{1}}{\isacharparenright}{\kern0pt}{\isacharbar}{\kern0pt}\ {\isasymand}\isanewline
\ \ \ \ scans\ {\isasymomega}\ k\ {\isacharparenleft}{\kern0pt}item\ {\isacharparenleft}{\kern0pt}bs{\isacharbang}{\kern0pt}{\isacharparenleft}{\kern0pt}k{\isacharminus}{\kern0pt}{\isadigit{1}}{\isacharparenright}{\kern0pt}{\isacharbang}{\kern0pt}pre{\isacharparenright}{\kern0pt}{\isacharparenright}{\kern0pt}\ {\isacharparenleft}{\kern0pt}item\ e{\isacharparenright}{\kern0pt}{\isachardoublequoteclose}\isanewline
\isanewline
\isacommand{definition}\isamarkupfalse%
\ sound{\isacharunderscore}{\kern0pt}prered{\isacharunderscore}{\kern0pt}ptr\ {\isacharcolon}{\kern0pt}{\isacharcolon}{\kern0pt}\ {\isachardoublequoteopen}{\isacharprime}{\kern0pt}a\ bins\ {\isasymRightarrow}\ nat\ {\isasymRightarrow}\ {\isacharprime}{\kern0pt}a\ entry\ {\isasymRightarrow}\ bool{\isachardoublequoteclose}\ \isakeyword{where}\isanewline
\ \ {\isachardoublequoteopen}sound{\isacharunderscore}{\kern0pt}prered{\isacharunderscore}{\kern0pt}ptr\ bs\ k\ e\ {\isasymequiv}\ {\isasymforall}p\ ps\ k{\isacharprime}{\kern0pt}\ pre\ red{\isachardot}{\kern0pt}\ pointer\ e\ {\isacharequal}{\kern0pt}\ PreRed\ p\ ps\ {\isasymand}\ {\isacharparenleft}{\kern0pt}k{\isacharprime}{\kern0pt}{\isacharcomma}{\kern0pt}\ pre{\isacharcomma}{\kern0pt}\ red{\isacharparenright}{\kern0pt}\ {\isasymin}\ set\ {\isacharparenleft}{\kern0pt}p{\isacharhash}{\kern0pt}ps{\isacharparenright}{\kern0pt}\ {\isasymlongrightarrow}\isanewline
\ \ \ \ k{\isacharprime}{\kern0pt}\ {\isacharless}{\kern0pt}\ k\ {\isasymand}\isanewline
\ \ \ \ pre\ {\isacharless}{\kern0pt}\ {\isacharbar}{\kern0pt}bs{\isacharbang}{\kern0pt}k{\isacharprime}{\kern0pt}{\isacharbar}{\kern0pt}\ {\isasymand}\ \isanewline
\ \ \ \ red\ {\isacharless}{\kern0pt}\ {\isacharbar}{\kern0pt}bs{\isacharbang}{\kern0pt}k{\isacharbar}{\kern0pt}\ {\isasymand}\isanewline
\ \ \ \ completes\ k\ {\isacharparenleft}{\kern0pt}item\ {\isacharparenleft}{\kern0pt}bs{\isacharbang}{\kern0pt}k{\isacharprime}{\kern0pt}{\isacharbang}{\kern0pt}pre{\isacharparenright}{\kern0pt}{\isacharparenright}{\kern0pt}\ {\isacharparenleft}{\kern0pt}item\ e{\isacharparenright}{\kern0pt}\ {\isacharparenleft}{\kern0pt}item\ {\isacharparenleft}{\kern0pt}bs{\isacharbang}{\kern0pt}k{\isacharbang}{\kern0pt}red{\isacharparenright}{\kern0pt}{\isacharparenright}{\kern0pt}{\isachardoublequoteclose}\isanewline
\isanewline
\isacommand{definition}\isamarkupfalse%
\ sound{\isacharunderscore}{\kern0pt}ptrs\ {\isacharcolon}{\kern0pt}{\isacharcolon}{\kern0pt}\ {\isachardoublequoteopen}{\isacharprime}{\kern0pt}a\ sentential\ {\isasymRightarrow}\ {\isacharprime}{\kern0pt}a\ bins\ {\isasymRightarrow}\ bool{\isachardoublequoteclose}\ \isakeyword{where}\isanewline
\ \ {\isachardoublequoteopen}sound{\isacharunderscore}{\kern0pt}ptrs\ {\isasymomega}\ bs\ {\isasymequiv}\ {\isasymforall}k\ {\isacharless}{\kern0pt}\ {\isacharbar}{\kern0pt}bs{\isacharbar}{\kern0pt}{\isachardot}{\kern0pt}\ {\isasymforall}e\ {\isasymin}\ set\ {\isacharparenleft}{\kern0pt}bs{\isacharbang}{\kern0pt}k{\isacharparenright}{\kern0pt}{\isachardot}{\kern0pt}\isanewline
\ \ \ \ sound{\isacharunderscore}{\kern0pt}null{\isacharunderscore}{\kern0pt}ptr\ e\ {\isasymand}\isanewline
\ \ \ \ sound{\isacharunderscore}{\kern0pt}pre{\isacharunderscore}{\kern0pt}ptr\ {\isasymomega}\ bs\ k\ e\ {\isasymand}\isanewline
\ \ \ \ sound{\isacharunderscore}{\kern0pt}prered{\isacharunderscore}{\kern0pt}ptr\ bs\ k\ e{\isachardoublequoteclose}\isanewline
\isanewline
\isacommand{definition}\isamarkupfalse%
\ mono{\isacharunderscore}{\kern0pt}red{\isacharunderscore}{\kern0pt}ptr\ {\isacharcolon}{\kern0pt}{\isacharcolon}{\kern0pt}\ {\isachardoublequoteopen}{\isacharprime}{\kern0pt}a\ bins\ {\isasymRightarrow}\ bool{\isachardoublequoteclose}\ \isakeyword{where}\isanewline
\ \ {\isachardoublequoteopen}mono{\isacharunderscore}{\kern0pt}red{\isacharunderscore}{\kern0pt}ptr\ bs\ {\isasymequiv}\ {\isasymforall}k\ {\isacharless}{\kern0pt}\ {\isacharbar}{\kern0pt}bs{\isacharbar}{\kern0pt}{\isachardot}{\kern0pt}\ {\isasymforall}i\ {\isacharless}{\kern0pt}\ {\isacharbar}{\kern0pt}bs{\isacharbang}{\kern0pt}k{\isacharbar}{\kern0pt}{\isachardot}{\kern0pt}\isanewline
\ \ \ \ {\isasymforall}k{\isacharprime}{\kern0pt}\ pre\ red\ ps{\isachardot}{\kern0pt}\ pointer\ {\isacharparenleft}{\kern0pt}bs{\isacharbang}{\kern0pt}k{\isacharbang}{\kern0pt}i{\isacharparenright}{\kern0pt}\ {\isacharequal}{\kern0pt}\ PreRed\ {\isacharparenleft}{\kern0pt}k{\isacharprime}{\kern0pt}{\isacharcomma}{\kern0pt}\ pre{\isacharcomma}{\kern0pt}\ red{\isacharparenright}{\kern0pt}\ ps\ {\isasymlongrightarrow}\ red\ {\isacharless}{\kern0pt}\ i{\isachardoublequoteclose}\isanewline
\isanewline
\isacommand{lemma}\isamarkupfalse%
\ sound{\isacharunderscore}{\kern0pt}mono{\isacharunderscore}{\kern0pt}ptrs{\isacharunderscore}{\kern0pt}bin{\isacharunderscore}{\kern0pt}upd{\isacharcolon}{\kern0pt}\isanewline
\ \ \isakeyword{assumes}\ {\isachardoublequoteopen}k\ {\isacharless}{\kern0pt}\ {\isacharbar}{\kern0pt}bs{\isacharbar}{\kern0pt}{\isachardoublequoteclose}\isanewline
\ \ \isakeyword{assumes}\ {\isachardoublequoteopen}distinct\ {\isacharparenleft}{\kern0pt}items\ {\isacharparenleft}{\kern0pt}bs{\isacharbang}{\kern0pt}k{\isacharparenright}{\kern0pt}{\isacharparenright}{\kern0pt}{\isachardoublequoteclose}\isanewline
\ \ \isakeyword{assumes}\ {\isachardoublequoteopen}sound{\isacharunderscore}{\kern0pt}ptrs\ {\isasymomega}\ bs{\isachardoublequoteclose}\isanewline
\ \ \isakeyword{assumes}\ {\isachardoublequoteopen}sound{\isacharunderscore}{\kern0pt}null{\isacharunderscore}{\kern0pt}ptr\ e{\isachardoublequoteclose}\isanewline
\ \ \isakeyword{assumes}\ {\isachardoublequoteopen}sound{\isacharunderscore}{\kern0pt}pre{\isacharunderscore}{\kern0pt}ptr\ {\isasymomega}\ bs\ k\ e{\isachardoublequoteclose}\isanewline
\ \ \isakeyword{assumes}\ {\isachardoublequoteopen}sound{\isacharunderscore}{\kern0pt}prered{\isacharunderscore}{\kern0pt}ptr\ bs\ k\ e{\isachardoublequoteclose}\isanewline
\ \ \isakeyword{assumes}\ {\isachardoublequoteopen}mono{\isacharunderscore}{\kern0pt}red{\isacharunderscore}{\kern0pt}ptr\ bs{\isachardoublequoteclose}\isanewline
\ \ \isakeyword{assumes}\ {\isachardoublequoteopen}{\isasymforall}k{\isacharprime}{\kern0pt}\ pre\ red\ ps{\isachardot}{\kern0pt}\ pointer\ e\ {\isacharequal}{\kern0pt}\ PreRed\ {\isacharparenleft}{\kern0pt}k{\isacharprime}{\kern0pt}{\isacharcomma}{\kern0pt}\ pre{\isacharcomma}{\kern0pt}\ red{\isacharparenright}{\kern0pt}\ ps\ {\isasymlongrightarrow}\ red\ {\isacharless}{\kern0pt}\ {\isacharbar}{\kern0pt}bs{\isacharbang}{\kern0pt}k{\isacharbar}{\kern0pt}{\isachardoublequoteclose}\isanewline
\ \ \isakeyword{shows}\ {\isachardoublequoteopen}sound{\isacharunderscore}{\kern0pt}ptrs\ {\isasymomega}\ {\isacharparenleft}{\kern0pt}bs{\isacharbrackleft}{\kern0pt}k\ {\isacharcolon}{\kern0pt}{\isacharequal}{\kern0pt}\ bin{\isacharunderscore}{\kern0pt}upd\ e\ {\isacharparenleft}{\kern0pt}bs{\isacharbang}{\kern0pt}k{\isacharparenright}{\kern0pt}{\isacharbrackright}{\kern0pt}{\isacharparenright}{\kern0pt}{\isachardoublequoteclose}%
\isadelimproof
%
\endisadelimproof
%
\isatagproof
%
\endisatagproof
{\isafoldproof}%
%
\isadelimproof
%
\endisadelimproof
%
\begin{isamarkuptext}%
%
\end{isamarkuptext}\isamarkuptrue%
\isacommand{lemma}\isamarkupfalse%
\ sound{\isacharunderscore}{\kern0pt}mono{\isacharunderscore}{\kern0pt}ptrs{\isacharunderscore}{\kern0pt}bin{\isacharunderscore}{\kern0pt}upds{\isacharcolon}{\kern0pt}\isanewline
\ \ \isakeyword{assumes}\ {\isachardoublequoteopen}k\ {\isacharless}{\kern0pt}\ {\isacharbar}{\kern0pt}bs{\isacharbar}{\kern0pt}{\isachardoublequoteclose}\isanewline
\ \ \isakeyword{assumes}\ {\isachardoublequoteopen}distinct\ {\isacharparenleft}{\kern0pt}items\ {\isacharparenleft}{\kern0pt}bs{\isacharbang}{\kern0pt}k{\isacharparenright}{\kern0pt}{\isacharparenright}{\kern0pt}{\isachardoublequoteclose}\isanewline
\ \ \isakeyword{assumes}\ {\isachardoublequoteopen}distinct\ {\isacharparenleft}{\kern0pt}items\ es{\isacharparenright}{\kern0pt}{\isachardoublequoteclose}\isanewline
\ \ \isakeyword{assumes}\ {\isachardoublequoteopen}sound{\isacharunderscore}{\kern0pt}ptrs\ inp\ bs{\isachardoublequoteclose}\isanewline
\ \ \isakeyword{assumes}\ {\isachardoublequoteopen}{\isasymforall}e\ {\isasymin}\ set\ es{\isachardot}{\kern0pt}\ sound{\isacharunderscore}{\kern0pt}null{\isacharunderscore}{\kern0pt}ptr\ e\ {\isasymand}\ sound{\isacharunderscore}{\kern0pt}pre{\isacharunderscore}{\kern0pt}ptr\ inp\ bs\ k\ e\ {\isasymand}\ sound{\isacharunderscore}{\kern0pt}prered{\isacharunderscore}{\kern0pt}ptr\ bs\ k\ e{\isachardoublequoteclose}\isanewline
\ \ \isakeyword{assumes}\ {\isachardoublequoteopen}mono{\isacharunderscore}{\kern0pt}red{\isacharunderscore}{\kern0pt}ptr\ bs{\isachardoublequoteclose}\isanewline
\ \ \isakeyword{assumes}\ {\isachardoublequoteopen}{\isasymforall}e\ {\isasymin}\ set\ es{\isachardot}{\kern0pt}\ {\isasymforall}k{\isacharprime}{\kern0pt}\ pre\ red\ ps{\isachardot}{\kern0pt}\ pointer\ e\ {\isacharequal}{\kern0pt}\ PreRed\ {\isacharparenleft}{\kern0pt}k{\isacharprime}{\kern0pt}{\isacharcomma}{\kern0pt}\ pre{\isacharcomma}{\kern0pt}\ red{\isacharparenright}{\kern0pt}\ ps\ {\isasymlongrightarrow}\ red\ {\isacharless}{\kern0pt}\ {\isacharbar}{\kern0pt}bs{\isacharbang}{\kern0pt}k{\isacharbar}{\kern0pt}{\isachardoublequoteclose}\isanewline
\ \ \isakeyword{shows}\ {\isachardoublequoteopen}sound{\isacharunderscore}{\kern0pt}ptrs\ inp\ {\isacharparenleft}{\kern0pt}bs{\isacharbrackleft}{\kern0pt}k\ {\isacharcolon}{\kern0pt}{\isacharequal}{\kern0pt}\ bin{\isacharunderscore}{\kern0pt}upds\ es\ {\isacharparenleft}{\kern0pt}bs{\isacharbang}{\kern0pt}k{\isacharparenright}{\kern0pt}{\isacharbrackright}{\kern0pt}{\isacharparenright}{\kern0pt}\ {\isasymand}\ mono{\isacharunderscore}{\kern0pt}red{\isacharunderscore}{\kern0pt}ptr\ {\isacharparenleft}{\kern0pt}bs{\isacharbrackleft}{\kern0pt}k\ {\isacharcolon}{\kern0pt}{\isacharequal}{\kern0pt}\ bin{\isacharunderscore}{\kern0pt}upds\ es\ {\isacharparenleft}{\kern0pt}bs{\isacharbang}{\kern0pt}k{\isacharparenright}{\kern0pt}{\isacharbrackright}{\kern0pt}{\isacharparenright}{\kern0pt}{\isachardoublequoteclose}%
\isadelimproof
%
\endisadelimproof
%
\isatagproof
%
\endisatagproof
{\isafoldproof}%
%
\isadelimproof
%
\endisadelimproof
%
\begin{isamarkuptext}%
%
\end{isamarkuptext}\isamarkuptrue%
\isacommand{lemma}\isamarkupfalse%
\ sound{\isacharunderscore}{\kern0pt}mono{\isacharunderscore}{\kern0pt}ptrs{\isacharunderscore}{\kern0pt}Earley{\isacharunderscore}{\kern0pt}bin{\isacharunderscore}{\kern0pt}list{\isacharprime}{\kern0pt}{\isacharcolon}{\kern0pt}\ %
\isamarkupcmt{Detailed%
}\isanewline
\ \ \isakeyword{assumes}\ {\isachardoublequoteopen}{\isacharparenleft}{\kern0pt}k{\isacharcomma}{\kern0pt}\ {\isasymG}{\isacharcomma}{\kern0pt}\ {\isasymomega}{\isacharcomma}{\kern0pt}\ bs{\isacharparenright}{\kern0pt}\ {\isasymin}\ wf{\isacharunderscore}{\kern0pt}earley{\isacharunderscore}{\kern0pt}input{\isachardoublequoteclose}\isanewline
\ \ \isakeyword{assumes}\ {\isachardoublequoteopen}nonempty{\isacharunderscore}{\kern0pt}derives\ {\isasymG}{\isachardoublequoteclose}\isanewline
\ \ \isakeyword{assumes}\ {\isachardoublequoteopen}sound{\isacharunderscore}{\kern0pt}items\ {\isasymG}\ {\isasymomega}\ {\isacharparenleft}{\kern0pt}bins\ bs{\isacharparenright}{\kern0pt}{\isachardoublequoteclose}\isanewline
\ \ \isakeyword{assumes}\ {\isachardoublequoteopen}sound{\isacharunderscore}{\kern0pt}ptrs\ {\isasymomega}\ bs{\isachardoublequoteclose}\ \isanewline
\ \ \isakeyword{assumes}\ {\isachardoublequoteopen}mono{\isacharunderscore}{\kern0pt}red{\isacharunderscore}{\kern0pt}ptr\ bs{\isachardoublequoteclose}\isanewline
\ \ \isakeyword{shows}\ {\isachardoublequoteopen}sound{\isacharunderscore}{\kern0pt}ptrs\ {\isasymomega}\ {\isacharparenleft}{\kern0pt}Earley{\isacharunderscore}{\kern0pt}bin{\isacharunderscore}{\kern0pt}list{\isacharprime}{\kern0pt}\ k\ {\isasymG}\ {\isasymomega}\ bs\ i{\isacharparenright}{\kern0pt}\ {\isasymand}\ mono{\isacharunderscore}{\kern0pt}red{\isacharunderscore}{\kern0pt}ptr\ {\isacharparenleft}{\kern0pt}Earley{\isacharunderscore}{\kern0pt}bin{\isacharunderscore}{\kern0pt}list{\isacharprime}{\kern0pt}\ k\ {\isasymG}\ {\isasymomega}\ bs\ i{\isacharparenright}{\kern0pt}{\isachardoublequoteclose}%
\isadelimproof
%
\endisadelimproof
%
\isatagproof
%
\endisatagproof
{\isafoldproof}%
%
\isadelimproof
%
\endisadelimproof
%
\begin{isamarkuptext}%
%
\end{isamarkuptext}\isamarkuptrue%
\isacommand{lemma}\isamarkupfalse%
\ sound{\isacharunderscore}{\kern0pt}mono{\isacharunderscore}{\kern0pt}ptrs{\isacharunderscore}{\kern0pt}Earley{\isacharunderscore}{\kern0pt}bin{\isacharunderscore}{\kern0pt}list{\isacharcolon}{\kern0pt}\isanewline
\ \ \isakeyword{assumes}\ {\isachardoublequoteopen}{\isacharparenleft}{\kern0pt}k{\isacharcomma}{\kern0pt}\ {\isasymG}{\isacharcomma}{\kern0pt}\ {\isasymomega}{\isacharcomma}{\kern0pt}\ bs{\isacharparenright}{\kern0pt}\ {\isasymin}\ wf{\isacharunderscore}{\kern0pt}earley{\isacharunderscore}{\kern0pt}input{\isachardoublequoteclose}\isanewline
\ \ \isakeyword{assumes}\ {\isachardoublequoteopen}nonempty{\isacharunderscore}{\kern0pt}derives\ {\isasymG}{\isachardoublequoteclose}\isanewline
\ \ \isakeyword{assumes}\ {\isachardoublequoteopen}sound{\isacharunderscore}{\kern0pt}items\ {\isasymG}\ {\isasymomega}\ {\isacharparenleft}{\kern0pt}bins\ bs{\isacharparenright}{\kern0pt}{\isachardoublequoteclose}\isanewline
\ \ \isakeyword{assumes}\ {\isachardoublequoteopen}sound{\isacharunderscore}{\kern0pt}ptrs\ {\isasymomega}\ bs{\isachardoublequoteclose}\isanewline
\ \ \isakeyword{assumes}\ {\isachardoublequoteopen}mono{\isacharunderscore}{\kern0pt}red{\isacharunderscore}{\kern0pt}ptr\ bs{\isachardoublequoteclose}\isanewline
\ \ \isakeyword{shows}\ {\isachardoublequoteopen}sound{\isacharunderscore}{\kern0pt}ptrs\ {\isasymomega}\ {\isacharparenleft}{\kern0pt}Earley{\isacharunderscore}{\kern0pt}bin{\isacharunderscore}{\kern0pt}list\ k\ {\isasymG}\ {\isasymomega}\ bs{\isacharparenright}{\kern0pt}\ {\isasymand}\ mono{\isacharunderscore}{\kern0pt}red{\isacharunderscore}{\kern0pt}ptr\ {\isacharparenleft}{\kern0pt}Earley{\isacharunderscore}{\kern0pt}bin{\isacharunderscore}{\kern0pt}list\ k\ {\isasymG}\ {\isasymomega}\ bs{\isacharparenright}{\kern0pt}{\isachardoublequoteclose}%
\isadelimproof
%
\endisadelimproof
%
\isatagproof
%
\endisatagproof
{\isafoldproof}%
%
\isadelimproof
%
\endisadelimproof
%
\begin{isamarkuptext}%
%
\end{isamarkuptext}\isamarkuptrue%
\isacommand{lemma}\isamarkupfalse%
\ sound{\isacharunderscore}{\kern0pt}mono{\isacharunderscore}{\kern0pt}ptrs{\isacharunderscore}{\kern0pt}Init{\isacharunderscore}{\kern0pt}list{\isacharcolon}{\kern0pt}\isanewline
\ \ \isakeyword{shows}\ {\isachardoublequoteopen}sound{\isacharunderscore}{\kern0pt}ptrs\ {\isasymomega}\ {\isacharparenleft}{\kern0pt}Init{\isacharunderscore}{\kern0pt}list\ {\isasymG}\ {\isasymomega}{\isacharparenright}{\kern0pt}\ {\isasymand}\ mono{\isacharunderscore}{\kern0pt}red{\isacharunderscore}{\kern0pt}ptr\ {\isacharparenleft}{\kern0pt}Init{\isacharunderscore}{\kern0pt}list\ {\isasymG}\ {\isasymomega}{\isacharparenright}{\kern0pt}{\isachardoublequoteclose}%
\isadelimproof
%
\endisadelimproof
%
\isatagproof
%
\endisatagproof
{\isafoldproof}%
%
\isadelimproof
%
\endisadelimproof
%
\begin{isamarkuptext}%
%
\end{isamarkuptext}\isamarkuptrue%
\isacommand{lemma}\isamarkupfalse%
\ sound{\isacharunderscore}{\kern0pt}mono{\isacharunderscore}{\kern0pt}ptrs{\isacharunderscore}{\kern0pt}Earley{\isacharunderscore}{\kern0pt}list{\isacharcolon}{\kern0pt}\isanewline
\ \ \isakeyword{assumes}\ {\isachardoublequoteopen}wf{\isacharunderscore}{\kern0pt}{\isasymG}\ {\isasymG}{\isachardoublequoteclose}\isanewline
\ \ \isakeyword{assumes}\ {\isachardoublequoteopen}nonempty{\isacharunderscore}{\kern0pt}derives\ {\isasymG}{\isachardoublequoteclose}\isanewline
\ \ \isakeyword{assumes}\ {\isachardoublequoteopen}k\ {\isasymle}\ {\isacharbar}{\kern0pt}{\isasymomega}{\isacharbar}{\kern0pt}{\isachardoublequoteclose}\isanewline
\ \ \isakeyword{shows}\ {\isachardoublequoteopen}sound{\isacharunderscore}{\kern0pt}ptrs\ {\isasymomega}\ {\isacharparenleft}{\kern0pt}Earley{\isacharunderscore}{\kern0pt}list\ k\ {\isasymG}\ {\isasymomega}{\isacharparenright}{\kern0pt}\ {\isasymand}\ mono{\isacharunderscore}{\kern0pt}red{\isacharunderscore}{\kern0pt}ptr\ {\isacharparenleft}{\kern0pt}Earley{\isacharunderscore}{\kern0pt}list\ k\ {\isasymG}\ {\isasymomega}{\isacharparenright}{\kern0pt}{\isachardoublequoteclose}%
\isadelimproof
%
\endisadelimproof
%
\isatagproof
%
\endisatagproof
{\isafoldproof}%
%
\isadelimproof
%
\endisadelimproof
%
\begin{isamarkuptext}%
%
\end{isamarkuptext}\isamarkuptrue%
\isacommand{lemma}\isamarkupfalse%
\ sound{\isacharunderscore}{\kern0pt}mono{\isacharunderscore}{\kern0pt}ptrs{\isacharunderscore}{\kern0pt}{\isasymE}arley{\isacharunderscore}{\kern0pt}list{\isacharcolon}{\kern0pt}\isanewline
\ \ \isakeyword{assumes}\ {\isachardoublequoteopen}wf{\isacharunderscore}{\kern0pt}{\isasymG}\ {\isasymG}{\isachardoublequoteclose}\isanewline
\ \ \isakeyword{assumes}\ {\isachardoublequoteopen}nonempty{\isacharunderscore}{\kern0pt}derives\ {\isasymG}{\isachardoublequoteclose}\isanewline
\ \ \isakeyword{shows}\ {\isachardoublequoteopen}sound{\isacharunderscore}{\kern0pt}ptrs\ {\isasymomega}\ {\isacharparenleft}{\kern0pt}{\isasymE}arley{\isacharunderscore}{\kern0pt}list\ {\isasymG}\ {\isasymomega}{\isacharparenright}{\kern0pt}\ {\isasymand}\ mono{\isacharunderscore}{\kern0pt}red{\isacharunderscore}{\kern0pt}ptr\ {\isacharparenleft}{\kern0pt}{\isasymE}arley{\isacharunderscore}{\kern0pt}list\ {\isasymG}\ {\isasymomega}{\isacharparenright}{\kern0pt}{\isachardoublequoteclose}%
\isadelimproof
%
\endisadelimproof
%
\isatagproof
%
\endisatagproof
{\isafoldproof}%
%
\isadelimproof
%
\endisadelimproof
%
\isadelimdocument
%
\endisadelimdocument
%
\isatagdocument
%
\isamarkupsubsection{The Parse Tree Algorithm%
}
\isamarkuptrue%
%
\endisatagdocument
{\isafolddocument}%
%
\isadelimdocument
%
\endisadelimdocument
\isacommand{partial{\isacharunderscore}{\kern0pt}function}\isamarkupfalse%
\ {\isacharparenleft}{\kern0pt}option{\isacharparenright}{\kern0pt}\ build{\isacharunderscore}{\kern0pt}tree{\isacharprime}{\kern0pt}\ {\isacharcolon}{\kern0pt}{\isacharcolon}{\kern0pt}\ {\isachardoublequoteopen}{\isacharprime}{\kern0pt}a\ bins\ {\isasymRightarrow}\ {\isacharprime}{\kern0pt}a\ sentential\ {\isasymRightarrow}\ nat\ {\isasymRightarrow}\ nat\ {\isasymRightarrow}\ {\isacharprime}{\kern0pt}a\ tree\ option{\isachardoublequoteclose}\ \isakeyword{where}\isanewline
\ \ {\isachardoublequoteopen}build{\isacharunderscore}{\kern0pt}tree{\isacharprime}{\kern0pt}\ bs\ {\isasymomega}\ k\ i\ {\isacharequal}{\kern0pt}\ {\isacharparenleft}{\kern0pt}\isanewline
\ \ \ \ let\ e\ {\isacharequal}{\kern0pt}\ bs{\isacharbang}{\kern0pt}k{\isacharbang}{\kern0pt}i\ in\ {\isacharparenleft}{\kern0pt}\isanewline
\ \ \ \ case\ pointer\ e\ of\isanewline
\ \ \ \ \ \ Null\ {\isasymRightarrow}\ Some\ {\isacharparenleft}{\kern0pt}Branch\ {\isacharparenleft}{\kern0pt}item{\isacharunderscore}{\kern0pt}rule{\isacharunderscore}{\kern0pt}head\ {\isacharparenleft}{\kern0pt}item\ e{\isacharparenright}{\kern0pt}{\isacharparenright}{\kern0pt}\ {\isacharbrackleft}{\kern0pt}{\isacharbrackright}{\kern0pt}{\isacharparenright}{\kern0pt}\isanewline
\ \ \ \ {\isacharbar}{\kern0pt}\ Pre\ pre\ {\isasymRightarrow}\ {\isacharparenleft}{\kern0pt}\isanewline
\ \ \ \ \ \ \ \ do\ {\isacharbraceleft}{\kern0pt}\isanewline
\ \ \ \ \ \ \ \ \ \ t\ {\isasymleftarrow}\ build{\isacharunderscore}{\kern0pt}tree{\isacharprime}{\kern0pt}\ bs\ {\isasymomega}\ {\isacharparenleft}{\kern0pt}k{\isacharminus}{\kern0pt}{\isadigit{1}}{\isacharparenright}{\kern0pt}\ pre{\isacharsemicolon}{\kern0pt}\isanewline
\ \ \ \ \ \ \ \ \ \ case\ t\ of\isanewline
\ \ \ \ \ \ \ \ \ \ \ \ Branch\ N\ ts\ {\isasymRightarrow}\ Some\ {\isacharparenleft}{\kern0pt}Branch\ N\ {\isacharparenleft}{\kern0pt}ts\ {\isacharat}{\kern0pt}\ {\isacharbrackleft}{\kern0pt}Leaf\ {\isacharparenleft}{\kern0pt}{\isasymomega}{\isacharbang}{\kern0pt}{\isacharparenleft}{\kern0pt}k{\isacharminus}{\kern0pt}{\isadigit{1}}{\isacharparenright}{\kern0pt}{\isacharparenright}{\kern0pt}{\isacharbrackright}{\kern0pt}{\isacharparenright}{\kern0pt}{\isacharparenright}{\kern0pt}\isanewline
\ \ \ \ \ \ \ \ \ \ {\isacharbar}{\kern0pt}\ {\isacharunderscore}{\kern0pt}\ {\isasymRightarrow}\ None\isanewline
\ \ \ \ \ \ \ \ {\isacharbraceright}{\kern0pt}{\isacharparenright}{\kern0pt}\isanewline
\ \ \ \ {\isacharbar}{\kern0pt}\ PreRed\ {\isacharparenleft}{\kern0pt}k{\isacharprime}{\kern0pt}{\isacharcomma}{\kern0pt}\ pre{\isacharcomma}{\kern0pt}\ red{\isacharparenright}{\kern0pt}\ {\isacharunderscore}{\kern0pt}\ {\isasymRightarrow}\ {\isacharparenleft}{\kern0pt}\isanewline
\ \ \ \ \ \ \ \ do\ {\isacharbraceleft}{\kern0pt}\isanewline
\ \ \ \ \ \ \ \ \ \ t\ {\isasymleftarrow}\ build{\isacharunderscore}{\kern0pt}tree{\isacharprime}{\kern0pt}\ bs\ {\isasymomega}\ k{\isacharprime}{\kern0pt}\ pre{\isacharsemicolon}{\kern0pt}\isanewline
\ \ \ \ \ \ \ \ \ \ case\ t\ of\isanewline
\ \ \ \ \ \ \ \ \ \ \ \ Branch\ N\ ts\ {\isasymRightarrow}\isanewline
\ \ \ \ \ \ \ \ \ \ \ \ \ \ do\ {\isacharbraceleft}{\kern0pt}\isanewline
\ \ \ \ \ \ \ \ \ \ \ \ \ \ \ \ t\ {\isasymleftarrow}\ build{\isacharunderscore}{\kern0pt}tree{\isacharprime}{\kern0pt}\ bs\ {\isasymomega}\ k\ red{\isacharsemicolon}{\kern0pt}\isanewline
\ \ \ \ \ \ \ \ \ \ \ \ \ \ \ \ Some\ {\isacharparenleft}{\kern0pt}Branch\ N\ {\isacharparenleft}{\kern0pt}ts\ {\isacharat}{\kern0pt}\ {\isacharbrackleft}{\kern0pt}t{\isacharbrackright}{\kern0pt}{\isacharparenright}{\kern0pt}{\isacharparenright}{\kern0pt}\isanewline
\ \ \ \ \ \ \ \ \ \ \ \ \ \ {\isacharbraceright}{\kern0pt}\isanewline
\ \ \ \ \ \ \ \ \ \ {\isacharbar}{\kern0pt}\ {\isacharunderscore}{\kern0pt}\ {\isasymRightarrow}\ None\isanewline
\ \ \ \ \ \ \ \ {\isacharbraceright}{\kern0pt}{\isacharparenright}{\kern0pt}\isanewline
\ \ {\isacharparenright}{\kern0pt}{\isacharparenright}{\kern0pt}{\isachardoublequoteclose}\isanewline
\isanewline
\isacommand{definition}\isamarkupfalse%
\ build{\isacharunderscore}{\kern0pt}tree\ {\isacharcolon}{\kern0pt}{\isacharcolon}{\kern0pt}\ {\isachardoublequoteopen}{\isacharprime}{\kern0pt}a\ cfg\ {\isasymRightarrow}\ {\isacharprime}{\kern0pt}a\ sentential\ {\isasymRightarrow}\ {\isacharprime}{\kern0pt}a\ bins\ {\isasymRightarrow}\ {\isacharprime}{\kern0pt}a\ tree\ option{\isachardoublequoteclose}\ \isakeyword{where}\isanewline
\ \ {\isachardoublequoteopen}build{\isacharunderscore}{\kern0pt}tree\ {\isasymG}\ {\isasymomega}\ bs\ {\isasymequiv}\isanewline
\ \ \ \ let\ k\ {\isacharequal}{\kern0pt}\ {\isacharbar}{\kern0pt}bs{\isacharbar}{\kern0pt}\ {\isacharminus}{\kern0pt}\ {\isadigit{1}}\ in\ {\isacharparenleft}{\kern0pt}\isanewline
\ \ \ \ case\ filter{\isacharunderscore}{\kern0pt}with{\isacharunderscore}{\kern0pt}index\ {\isacharparenleft}{\kern0pt}{\isasymlambda}x{\isachardot}{\kern0pt}\ is{\isacharunderscore}{\kern0pt}finished\ {\isasymG}\ {\isasymomega}\ x{\isacharparenright}{\kern0pt}\ {\isacharparenleft}{\kern0pt}items\ {\isacharparenleft}{\kern0pt}bs{\isacharbang}{\kern0pt}k{\isacharparenright}{\kern0pt}{\isacharparenright}{\kern0pt}\ of\isanewline
\ \ \ \ \ \ {\isacharbrackleft}{\kern0pt}{\isacharbrackright}{\kern0pt}\ {\isasymRightarrow}\ None\isanewline
\ \ \ \ {\isacharbar}{\kern0pt}\ {\isacharparenleft}{\kern0pt}{\isacharunderscore}{\kern0pt}{\isacharcomma}{\kern0pt}\ i{\isacharparenright}{\kern0pt}{\isacharhash}{\kern0pt}{\isacharunderscore}{\kern0pt}\ {\isasymRightarrow}\ build{\isacharunderscore}{\kern0pt}tree{\isacharprime}{\kern0pt}\ bs\ {\isasymomega}\ k\ i{\isacharparenright}{\kern0pt}{\isachardoublequoteclose}%
\isadelimdocument
%
\endisadelimdocument
%
\isatagdocument
%
\isamarkupsubsection{Termination%
}
\isamarkuptrue%
%
\endisatagdocument
{\isafolddocument}%
%
\isadelimdocument
%
\endisadelimdocument
\isacommand{fun}\isamarkupfalse%
\ build{\isacharunderscore}{\kern0pt}tree{\isacharprime}{\kern0pt}{\isacharunderscore}{\kern0pt}measure\ {\isacharcolon}{\kern0pt}{\isacharcolon}{\kern0pt}\ {\isachardoublequoteopen}{\isacharparenleft}{\kern0pt}{\isacharprime}{\kern0pt}a\ bins\ {\isasymtimes}\ {\isacharprime}{\kern0pt}a\ sentential\ {\isasymtimes}\ nat\ {\isasymtimes}\ nat{\isacharparenright}{\kern0pt}\ {\isasymRightarrow}\ nat{\isachardoublequoteclose}\ \isakeyword{where}\isanewline
\ \ {\isachardoublequoteopen}build{\isacharunderscore}{\kern0pt}tree{\isacharprime}{\kern0pt}{\isacharunderscore}{\kern0pt}measure\ {\isacharparenleft}{\kern0pt}bs{\isacharcomma}{\kern0pt}\ {\isasymomega}{\isacharcomma}{\kern0pt}\ k{\isacharcomma}{\kern0pt}\ i{\isacharparenright}{\kern0pt}\ {\isacharequal}{\kern0pt}\ foldl\ {\isacharparenleft}{\kern0pt}{\isacharplus}{\kern0pt}{\isacharparenright}{\kern0pt}\ {\isadigit{0}}\ {\isacharparenleft}{\kern0pt}map\ length\ {\isacharparenleft}{\kern0pt}take\ k\ bs{\isacharparenright}{\kern0pt}{\isacharparenright}{\kern0pt}\ {\isacharplus}{\kern0pt}\ i{\isachardoublequoteclose}\isanewline
\isanewline
\isacommand{definition}\isamarkupfalse%
\ wf{\isacharunderscore}{\kern0pt}tree{\isacharunderscore}{\kern0pt}input\ {\isacharcolon}{\kern0pt}{\isacharcolon}{\kern0pt}\ {\isachardoublequoteopen}{\isacharparenleft}{\kern0pt}{\isacharprime}{\kern0pt}a\ bins\ {\isasymtimes}\ {\isacharprime}{\kern0pt}a\ sentential\ {\isasymtimes}\ nat\ {\isasymtimes}\ nat{\isacharparenright}{\kern0pt}\ set{\isachardoublequoteclose}\ \isakeyword{where}\isanewline
\ \ {\isachardoublequoteopen}wf{\isacharunderscore}{\kern0pt}tree{\isacharunderscore}{\kern0pt}input\ {\isacharequal}{\kern0pt}\ {\isacharbraceleft}{\kern0pt}\isanewline
\ \ \ \ {\isacharparenleft}{\kern0pt}bs{\isacharcomma}{\kern0pt}\ {\isasymomega}{\isacharcomma}{\kern0pt}\ k{\isacharcomma}{\kern0pt}\ i{\isacharparenright}{\kern0pt}\ {\isacharbar}{\kern0pt}\ bs\ {\isasymomega}\ k\ i{\isachardot}{\kern0pt}\isanewline
\ \ \ \ \ \ sound{\isacharunderscore}{\kern0pt}ptrs\ {\isasymomega}\ bs\ {\isasymand}\isanewline
\ \ \ \ \ \ mono{\isacharunderscore}{\kern0pt}red{\isacharunderscore}{\kern0pt}ptr\ bs\ {\isasymand}\isanewline
\ \ \ \ \ \ k\ {\isacharless}{\kern0pt}\ {\isacharbar}{\kern0pt}bs{\isacharbar}{\kern0pt}\ {\isasymand}\isanewline
\ \ \ \ \ \ i\ {\isacharless}{\kern0pt}\ {\isacharbar}{\kern0pt}bs{\isacharbang}{\kern0pt}k{\isacharbar}{\kern0pt}\isanewline
\ \ {\isacharbraceright}{\kern0pt}{\isachardoublequoteclose}\isanewline
\isanewline
\isacommand{lemma}\isamarkupfalse%
\ build{\isacharunderscore}{\kern0pt}tree{\isacharprime}{\kern0pt}{\isacharunderscore}{\kern0pt}termination{\isacharcolon}{\kern0pt}\isanewline
\ \ \isakeyword{assumes}\ {\isachardoublequoteopen}{\isacharparenleft}{\kern0pt}bs{\isacharcomma}{\kern0pt}\ {\isasymomega}{\isacharcomma}{\kern0pt}\ k{\isacharcomma}{\kern0pt}\ i{\isacharparenright}{\kern0pt}\ {\isasymin}\ wf{\isacharunderscore}{\kern0pt}tree{\isacharunderscore}{\kern0pt}input{\isachardoublequoteclose}\isanewline
\ \ \isakeyword{shows}\ {\isachardoublequoteopen}{\isasymexists}N\ ts{\isachardot}{\kern0pt}\ build{\isacharunderscore}{\kern0pt}tree{\isacharprime}{\kern0pt}\ bs\ {\isasymomega}\ k\ i\ {\isacharequal}{\kern0pt}\ Some\ {\isacharparenleft}{\kern0pt}Branch\ N\ ts{\isacharparenright}{\kern0pt}{\isachardoublequoteclose}%
\isadelimproof
%
\endisadelimproof
%
\isatagproof
%
\endisatagproof
{\isafoldproof}%
%
\isadelimproof
%
\endisadelimproof
%
\begin{isamarkuptext}%
%
\end{isamarkuptext}\isamarkuptrue%
%
\isadelimdocument
%
\endisadelimdocument
%
\isatagdocument
%
\isamarkupsubsection{Soundness%
}
\isamarkuptrue%
%
\endisatagdocument
{\isafolddocument}%
%
\isadelimdocument
%
\endisadelimdocument
\isacommand{lemma}\isamarkupfalse%
\ wf{\isacharunderscore}{\kern0pt}item{\isacharunderscore}{\kern0pt}tree{\isacharunderscore}{\kern0pt}build{\isacharunderscore}{\kern0pt}tree{\isacharprime}{\kern0pt}{\isacharcolon}{\kern0pt}\isanewline
\ \ \isakeyword{assumes}\ {\isachardoublequoteopen}{\isacharparenleft}{\kern0pt}bs{\isacharcomma}{\kern0pt}\ {\isasymomega}{\isacharcomma}{\kern0pt}\ k{\isacharcomma}{\kern0pt}\ i{\isacharparenright}{\kern0pt}\ {\isasymin}\ wf{\isacharunderscore}{\kern0pt}tree{\isacharunderscore}{\kern0pt}input{\isachardoublequoteclose}\isanewline
\ \ \isakeyword{assumes}\ {\isachardoublequoteopen}wf{\isacharunderscore}{\kern0pt}bins\ {\isasymG}\ {\isasymomega}\ bs{\isachardoublequoteclose}\isanewline
\ \ \isakeyword{assumes}\ {\isachardoublequoteopen}k\ {\isacharless}{\kern0pt}\ {\isacharbar}{\kern0pt}bs{\isacharbar}{\kern0pt}{\isachardoublequoteclose}\isanewline
\ \ \isakeyword{assumes}\ {\isachardoublequoteopen}i\ {\isacharless}{\kern0pt}\ {\isacharbar}{\kern0pt}bs{\isacharbang}{\kern0pt}k{\isacharbar}{\kern0pt}{\isachardoublequoteclose}\isanewline
\ \ \isakeyword{assumes}\ {\isachardoublequoteopen}build{\isacharunderscore}{\kern0pt}tree{\isacharprime}{\kern0pt}\ bs\ {\isasymomega}\ k\ i\ {\isacharequal}{\kern0pt}\ Some\ t{\isachardoublequoteclose}\isanewline
\ \ \isakeyword{shows}\ {\isachardoublequoteopen}wf{\isacharunderscore}{\kern0pt}item{\isacharunderscore}{\kern0pt}tree\ {\isasymG}\ {\isacharparenleft}{\kern0pt}item\ {\isacharparenleft}{\kern0pt}bs{\isacharbang}{\kern0pt}k{\isacharbang}{\kern0pt}i{\isacharparenright}{\kern0pt}{\isacharparenright}{\kern0pt}\ t{\isachardoublequoteclose}%
\isadelimproof
%
\endisadelimproof
%
\isatagproof
%
\endisatagproof
{\isafoldproof}%
%
\isadelimproof
%
\endisadelimproof
%
\begin{isamarkuptext}%
%
\end{isamarkuptext}\isamarkuptrue%
\isacommand{lemma}\isamarkupfalse%
\ wf{\isacharunderscore}{\kern0pt}yield{\isacharunderscore}{\kern0pt}tree{\isacharunderscore}{\kern0pt}build{\isacharunderscore}{\kern0pt}tree{\isacharprime}{\kern0pt}{\isacharcolon}{\kern0pt}\isanewline
\ \ \isakeyword{assumes}\ {\isachardoublequoteopen}{\isacharparenleft}{\kern0pt}bs{\isacharcomma}{\kern0pt}\ {\isasymomega}{\isacharcomma}{\kern0pt}\ k{\isacharcomma}{\kern0pt}\ i{\isacharparenright}{\kern0pt}\ {\isasymin}\ wf{\isacharunderscore}{\kern0pt}tree{\isacharunderscore}{\kern0pt}input{\isachardoublequoteclose}\isanewline
\ \ \isakeyword{assumes}\ {\isachardoublequoteopen}wf{\isacharunderscore}{\kern0pt}bins\ {\isasymG}\ {\isasymomega}\ bs{\isachardoublequoteclose}\isanewline
\ \ \isakeyword{assumes}\ {\isachardoublequoteopen}k\ {\isacharless}{\kern0pt}\ {\isacharbar}{\kern0pt}bs{\isacharbar}{\kern0pt}{\isachardoublequoteclose}\isanewline
\ \ \isakeyword{assumes}\ {\isachardoublequoteopen}k\ {\isasymle}\ {\isacharbar}{\kern0pt}{\isasymomega}{\isacharbar}{\kern0pt}{\isachardoublequoteclose}\isanewline
\ \ \isakeyword{assumes}\ {\isachardoublequoteopen}i\ {\isacharless}{\kern0pt}\ {\isacharbar}{\kern0pt}bs{\isacharbang}{\kern0pt}k{\isacharbar}{\kern0pt}{\isachardoublequoteclose}\isanewline
\ \ \isakeyword{assumes}\ {\isachardoublequoteopen}build{\isacharunderscore}{\kern0pt}tree{\isacharprime}{\kern0pt}\ bs\ {\isasymomega}\ k\ i\ {\isacharequal}{\kern0pt}\ Some\ t{\isachardoublequoteclose}\isanewline
\ \ \isakeyword{shows}\ {\isachardoublequoteopen}wf{\isacharunderscore}{\kern0pt}yield{\isacharunderscore}{\kern0pt}tree\ {\isasymomega}\ {\isacharparenleft}{\kern0pt}item\ {\isacharparenleft}{\kern0pt}bs{\isacharbang}{\kern0pt}k{\isacharbang}{\kern0pt}i{\isacharparenright}{\kern0pt}{\isacharparenright}{\kern0pt}\ t{\isachardoublequoteclose}%
\isadelimproof
%
\endisadelimproof
%
\isatagproof
%
\endisatagproof
{\isafoldproof}%
%
\isadelimproof
%
\endisadelimproof
%
\begin{isamarkuptext}%
%
\end{isamarkuptext}\isamarkuptrue%
\isacommand{theorem}\isamarkupfalse%
\ wf{\isacharunderscore}{\kern0pt}rule{\isacharunderscore}{\kern0pt}root{\isacharunderscore}{\kern0pt}yield{\isacharunderscore}{\kern0pt}tree{\isacharunderscore}{\kern0pt}build{\isacharunderscore}{\kern0pt}tree{\isacharcolon}{\kern0pt}\isanewline
\ \ \isakeyword{assumes}\ {\isachardoublequoteopen}wf{\isacharunderscore}{\kern0pt}bins\ {\isasymG}\ {\isasymomega}\ bs{\isachardoublequoteclose}\isanewline
\ \ \isakeyword{assumes}\ {\isachardoublequoteopen}sound{\isacharunderscore}{\kern0pt}ptrs\ {\isasymomega}\ bs{\isachardoublequoteclose}\isanewline
\ \ \isakeyword{assumes}\ {\isachardoublequoteopen}mono{\isacharunderscore}{\kern0pt}red{\isacharunderscore}{\kern0pt}ptr\ bs{\isachardoublequoteclose}\isanewline
\ \ \isakeyword{assumes}\ {\isachardoublequoteopen}{\isacharbar}{\kern0pt}bs{\isacharbar}{\kern0pt}\ {\isacharequal}{\kern0pt}\ {\isacharbar}{\kern0pt}{\isasymomega}{\isacharbar}{\kern0pt}\ {\isacharplus}{\kern0pt}\ {\isadigit{1}}{\isachardoublequoteclose}\isanewline
\ \ \isakeyword{assumes}\ {\isachardoublequoteopen}build{\isacharunderscore}{\kern0pt}tree\ {\isasymG}\ {\isasymomega}\ bs\ {\isacharequal}{\kern0pt}\ Some\ t{\isachardoublequoteclose}\isanewline
\ \ \isakeyword{shows}\ {\isachardoublequoteopen}wf{\isacharunderscore}{\kern0pt}rule{\isacharunderscore}{\kern0pt}tree\ {\isasymG}\ t\ {\isasymand}\ root{\isacharunderscore}{\kern0pt}tree\ t\ {\isacharequal}{\kern0pt}\ {\isasymSS}\ {\isasymG}\ {\isasymand}\ yield{\isacharunderscore}{\kern0pt}tree\ t\ {\isacharequal}{\kern0pt}\ {\isasymomega}{\isachardoublequoteclose}%
\isadelimproof
%
\endisadelimproof
%
\isatagproof
%
\endisatagproof
{\isafoldproof}%
%
\isadelimproof
%
\endisadelimproof
%
\begin{isamarkuptext}%
%
\end{isamarkuptext}\isamarkuptrue%
\isacommand{corollary}\isamarkupfalse%
\ wf{\isacharunderscore}{\kern0pt}rule{\isacharunderscore}{\kern0pt}root{\isacharunderscore}{\kern0pt}yield{\isacharunderscore}{\kern0pt}tree{\isacharunderscore}{\kern0pt}build{\isacharunderscore}{\kern0pt}tree{\isacharunderscore}{\kern0pt}{\isasymE}arley{\isacharunderscore}{\kern0pt}list{\isacharcolon}{\kern0pt}\isanewline
\ \ \isakeyword{assumes}\ {\isachardoublequoteopen}wf{\isacharunderscore}{\kern0pt}{\isasymG}\ {\isasymG}{\isachardoublequoteclose}\isanewline
\ \ \isakeyword{assumes}\ {\isachardoublequoteopen}nonempty{\isacharunderscore}{\kern0pt}derives\ {\isasymG}{\isachardoublequoteclose}\isanewline
\ \ \isakeyword{assumes}\ {\isachardoublequoteopen}build{\isacharunderscore}{\kern0pt}tree\ {\isasymG}\ {\isasymomega}\ {\isacharparenleft}{\kern0pt}{\isasymE}arley{\isacharunderscore}{\kern0pt}list\ {\isasymG}\ {\isasymomega}{\isacharparenright}{\kern0pt}\ {\isacharequal}{\kern0pt}\ Some\ t{\isachardoublequoteclose}\isanewline
\ \ \isakeyword{shows}\ {\isachardoublequoteopen}wf{\isacharunderscore}{\kern0pt}rule{\isacharunderscore}{\kern0pt}tree\ {\isasymG}\ t\ {\isasymand}\ root{\isacharunderscore}{\kern0pt}tree\ t\ {\isacharequal}{\kern0pt}\ {\isasymSS}\ {\isasymG}\ {\isasymand}\ yield{\isacharunderscore}{\kern0pt}tree\ t\ {\isacharequal}{\kern0pt}\ {\isasymomega}{\isachardoublequoteclose}%
\isadelimproof
%
\endisadelimproof
%
\isatagproof
%
\endisatagproof
{\isafoldproof}%
%
\isadelimproof
%
\endisadelimproof
%
\begin{isamarkuptext}%
%
\end{isamarkuptext}\isamarkuptrue%
%
\isadelimdocument
%
\endisadelimdocument
%
\isatagdocument
%
\isamarkupsubsection{Completeness%
}
\isamarkuptrue%
%
\endisatagdocument
{\isafolddocument}%
%
\isadelimdocument
%
\endisadelimdocument
\isacommand{theorem}\isamarkupfalse%
\ correctness{\isacharunderscore}{\kern0pt}build{\isacharunderscore}{\kern0pt}tree{\isacharunderscore}{\kern0pt}{\isasymE}arley{\isacharunderscore}{\kern0pt}list{\isacharcolon}{\kern0pt}\isanewline
\ \ \isakeyword{assumes}\ {\isachardoublequoteopen}wf{\isacharunderscore}{\kern0pt}{\isasymG}\ {\isasymG}{\isachardoublequoteclose}\isanewline
\ \ \isakeyword{assumes}\ {\isachardoublequoteopen}is{\isacharunderscore}{\kern0pt}sentence\ {\isasymG}\ {\isasymomega}{\isachardoublequoteclose}\isanewline
\ \ \isakeyword{assumes}\ {\isachardoublequoteopen}nonempty{\isacharunderscore}{\kern0pt}derives\ {\isasymG}{\isachardoublequoteclose}\isanewline
\ \ \isakeyword{shows}\ {\isachardoublequoteopen}{\isacharparenleft}{\kern0pt}{\isasymexists}t{\isachardot}{\kern0pt}\ build{\isacharunderscore}{\kern0pt}tree\ {\isasymG}\ {\isasymomega}\ {\isacharparenleft}{\kern0pt}{\isasymE}arley{\isacharunderscore}{\kern0pt}list\ {\isasymG}\ {\isasymomega}{\isacharparenright}{\kern0pt}\ {\isacharequal}{\kern0pt}\ Some\ t{\isacharparenright}{\kern0pt}\ {\isasymlongleftrightarrow}\ {\isasymG}\ {\isasymturnstile}\ {\isacharbrackleft}{\kern0pt}{\isasymSS}\ {\isasymG}{\isacharbrackright}{\kern0pt}\ {\isasymRightarrow}\isactrlsup {\isacharasterisk}{\kern0pt}\ {\isasymomega}{\isachardoublequoteclose}%
\isadelimproof
%
\endisadelimproof
%
\isatagproof
%
\endisatagproof
{\isafoldproof}%
%
\isadelimproof
%
\endisadelimproof
%
\isadelimdocument
%
\endisadelimdocument
%
\isatagdocument
%
\isamarkupsection{A Parse Forest%
}
\isamarkuptrue%
%
\endisatagdocument
{\isafolddocument}%
%
\isadelimdocument
%
\endisadelimdocument
\isacommand{datatype}\isamarkupfalse%
\ {\isacharprime}{\kern0pt}a\ forest\ {\isacharequal}{\kern0pt}\isanewline
\ \ FLeaf\ {\isacharprime}{\kern0pt}a\isanewline
\ \ {\isacharbar}{\kern0pt}\ FBranch\ {\isacharprime}{\kern0pt}a\ {\isachardoublequoteopen}{\isacharprime}{\kern0pt}a\ forest\ list\ list{\isachardoublequoteclose}\isanewline
\isanewline
\isacommand{fun}\isamarkupfalse%
\ combinations\ {\isacharcolon}{\kern0pt}{\isacharcolon}{\kern0pt}\ {\isachardoublequoteopen}{\isacharprime}{\kern0pt}a\ list\ list\ {\isasymRightarrow}\ {\isacharprime}{\kern0pt}a\ list\ list{\isachardoublequoteclose}\ \isakeyword{where}\isanewline
\ \ {\isachardoublequoteopen}combinations\ {\isacharbrackleft}{\kern0pt}{\isacharbrackright}{\kern0pt}\ {\isacharequal}{\kern0pt}\ {\isacharbrackleft}{\kern0pt}{\isacharbrackleft}{\kern0pt}{\isacharbrackright}{\kern0pt}{\isacharbrackright}{\kern0pt}{\isachardoublequoteclose}\isanewline
{\isacharbar}{\kern0pt}\ {\isachardoublequoteopen}combinations\ {\isacharparenleft}{\kern0pt}xs{\isacharhash}{\kern0pt}xss{\isacharparenright}{\kern0pt}\ {\isacharequal}{\kern0pt}\ {\isacharbrackleft}{\kern0pt}\ x{\isacharhash}{\kern0pt}cs\ {\isachardot}{\kern0pt}\ x\ {\isacharless}{\kern0pt}{\isacharminus}{\kern0pt}\ xs{\isacharcomma}{\kern0pt}\ cs\ {\isacharless}{\kern0pt}{\isacharminus}{\kern0pt}\ combinations\ xss\ {\isacharbrackright}{\kern0pt}{\isachardoublequoteclose}\isanewline
\isanewline
\isacommand{fun}\isamarkupfalse%
\ trees\ {\isacharcolon}{\kern0pt}{\isacharcolon}{\kern0pt}\ {\isachardoublequoteopen}{\isacharprime}{\kern0pt}a\ forest\ {\isasymRightarrow}\ {\isacharprime}{\kern0pt}a\ tree\ list{\isachardoublequoteclose}\ \isakeyword{where}\isanewline
\ \ {\isachardoublequoteopen}trees\ {\isacharparenleft}{\kern0pt}FLeaf\ a{\isacharparenright}{\kern0pt}\ {\isacharequal}{\kern0pt}\ {\isacharbrackleft}{\kern0pt}Leaf\ a{\isacharbrackright}{\kern0pt}{\isachardoublequoteclose}\isanewline
{\isacharbar}{\kern0pt}\ {\isachardoublequoteopen}trees\ {\isacharparenleft}{\kern0pt}FBranch\ N\ fss{\isacharparenright}{\kern0pt}\ {\isacharequal}{\kern0pt}\ {\isacharparenleft}{\kern0pt}\isanewline
\ \ \ \ let\ tss\ {\isacharequal}{\kern0pt}\ {\isacharparenleft}{\kern0pt}map\ {\isacharparenleft}{\kern0pt}{\isasymlambda}fs{\isachardot}{\kern0pt}\ concat\ {\isacharparenleft}{\kern0pt}map\ {\isacharparenleft}{\kern0pt}{\isasymlambda}f{\isachardot}{\kern0pt}\ trees\ f{\isacharparenright}{\kern0pt}\ fs{\isacharparenright}{\kern0pt}{\isacharparenright}{\kern0pt}\ fss{\isacharparenright}{\kern0pt}\ in\isanewline
\ \ \ \ map\ {\isacharparenleft}{\kern0pt}{\isasymlambda}ts{\isachardot}{\kern0pt}\ Branch\ N\ ts{\isacharparenright}{\kern0pt}\ {\isacharparenleft}{\kern0pt}combinations\ tss{\isacharparenright}{\kern0pt}\isanewline
\ \ {\isacharparenright}{\kern0pt}{\isachardoublequoteclose}%
\isadelimdocument
%
\endisadelimdocument
%
\isatagdocument
%
\isamarkupsubsection{The Parse Forest Algorithm%
}
\isamarkuptrue%
%
\endisatagdocument
{\isafolddocument}%
%
\isadelimdocument
%
\endisadelimdocument
\isacommand{fun}\isamarkupfalse%
\ insert{\isacharunderscore}{\kern0pt}group\ {\isacharcolon}{\kern0pt}{\isacharcolon}{\kern0pt}\ {\isachardoublequoteopen}{\isacharparenleft}{\kern0pt}{\isacharprime}{\kern0pt}a\ {\isasymRightarrow}\ {\isacharprime}{\kern0pt}k{\isacharparenright}{\kern0pt}\ {\isasymRightarrow}\ {\isacharparenleft}{\kern0pt}{\isacharprime}{\kern0pt}a\ {\isasymRightarrow}\ {\isacharprime}{\kern0pt}v{\isacharparenright}{\kern0pt}\ {\isasymRightarrow}\ {\isacharprime}{\kern0pt}a\ {\isasymRightarrow}\ {\isacharparenleft}{\kern0pt}{\isacharprime}{\kern0pt}k\ {\isasymtimes}\ {\isacharprime}{\kern0pt}v\ list{\isacharparenright}{\kern0pt}\ list\ {\isasymRightarrow}\ {\isacharparenleft}{\kern0pt}{\isacharprime}{\kern0pt}k\ {\isasymtimes}\ {\isacharprime}{\kern0pt}v\ list{\isacharparenright}{\kern0pt}\ list{\isachardoublequoteclose}\ \isakeyword{where}\isanewline
\ \ {\isachardoublequoteopen}insert{\isacharunderscore}{\kern0pt}group\ K\ V\ a\ {\isacharbrackleft}{\kern0pt}{\isacharbrackright}{\kern0pt}\ {\isacharequal}{\kern0pt}\ {\isacharbrackleft}{\kern0pt}{\isacharparenleft}{\kern0pt}K\ a{\isacharcomma}{\kern0pt}\ {\isacharbrackleft}{\kern0pt}V\ a{\isacharbrackright}{\kern0pt}{\isacharparenright}{\kern0pt}{\isacharbrackright}{\kern0pt}{\isachardoublequoteclose}\isanewline
{\isacharbar}{\kern0pt}\ {\isachardoublequoteopen}insert{\isacharunderscore}{\kern0pt}group\ K\ V\ a\ {\isacharparenleft}{\kern0pt}{\isacharparenleft}{\kern0pt}k{\isacharcomma}{\kern0pt}\ vs{\isacharparenright}{\kern0pt}{\isacharhash}{\kern0pt}xs{\isacharparenright}{\kern0pt}\ {\isacharequal}{\kern0pt}\ {\isacharparenleft}{\kern0pt}\isanewline
\ \ \ \ if\ K\ a\ {\isacharequal}{\kern0pt}\ k\ then\ {\isacharparenleft}{\kern0pt}k{\isacharcomma}{\kern0pt}\ V\ a\ {\isacharhash}{\kern0pt}\ vs{\isacharparenright}{\kern0pt}\ {\isacharhash}{\kern0pt}\ xs\isanewline
\ \ \ \ else\ {\isacharparenleft}{\kern0pt}k{\isacharcomma}{\kern0pt}\ vs{\isacharparenright}{\kern0pt}\ {\isacharhash}{\kern0pt}\ insert{\isacharunderscore}{\kern0pt}group\ K\ V\ a\ xs\ \ \isanewline
\ \ {\isacharparenright}{\kern0pt}{\isachardoublequoteclose}\isanewline
\isanewline
\isacommand{fun}\isamarkupfalse%
\ group{\isacharunderscore}{\kern0pt}by\ {\isacharcolon}{\kern0pt}{\isacharcolon}{\kern0pt}\ {\isachardoublequoteopen}{\isacharparenleft}{\kern0pt}{\isacharprime}{\kern0pt}a\ {\isasymRightarrow}\ {\isacharprime}{\kern0pt}k{\isacharparenright}{\kern0pt}\ {\isasymRightarrow}\ {\isacharparenleft}{\kern0pt}{\isacharprime}{\kern0pt}a\ {\isasymRightarrow}\ {\isacharprime}{\kern0pt}v{\isacharparenright}{\kern0pt}\ {\isasymRightarrow}\ {\isacharprime}{\kern0pt}a\ list\ {\isasymRightarrow}\ {\isacharparenleft}{\kern0pt}{\isacharprime}{\kern0pt}k\ {\isasymtimes}\ {\isacharprime}{\kern0pt}v\ list{\isacharparenright}{\kern0pt}\ list{\isachardoublequoteclose}\ \isakeyword{where}\isanewline
\ \ {\isachardoublequoteopen}group{\isacharunderscore}{\kern0pt}by\ K\ V\ {\isacharbrackleft}{\kern0pt}{\isacharbrackright}{\kern0pt}\ {\isacharequal}{\kern0pt}\ {\isacharbrackleft}{\kern0pt}{\isacharbrackright}{\kern0pt}{\isachardoublequoteclose}\isanewline
{\isacharbar}{\kern0pt}\ {\isachardoublequoteopen}group{\isacharunderscore}{\kern0pt}by\ K\ V\ {\isacharparenleft}{\kern0pt}x{\isacharhash}{\kern0pt}xs{\isacharparenright}{\kern0pt}\ {\isacharequal}{\kern0pt}\ insert{\isacharunderscore}{\kern0pt}group\ K\ V\ x\ {\isacharparenleft}{\kern0pt}group{\isacharunderscore}{\kern0pt}by\ K\ V\ xs{\isacharparenright}{\kern0pt}{\isachardoublequoteclose}\isanewline
%
\isadelimproof
%
\endisadelimproof
%
\isatagproof
%
\endisatagproof
{\isafoldproof}%
%
\isadelimproof
\isanewline
%
\endisadelimproof
\isacommand{partial{\isacharunderscore}{\kern0pt}function}\isamarkupfalse%
\ {\isacharparenleft}{\kern0pt}option{\isacharparenright}{\kern0pt}\ build{\isacharunderscore}{\kern0pt}trees{\isacharprime}{\kern0pt}\ {\isacharcolon}{\kern0pt}{\isacharcolon}{\kern0pt}\ {\isachardoublequoteopen}{\isacharprime}{\kern0pt}a\ bins\ {\isasymRightarrow}\ {\isacharprime}{\kern0pt}a\ sentential\ {\isasymRightarrow}\ nat\ {\isasymRightarrow}\ nat\ {\isasymRightarrow}\ nat\ set\ {\isasymRightarrow}\ {\isacharprime}{\kern0pt}a\ forest\ list\ option{\isachardoublequoteclose}\ \isakeyword{where}\isanewline
\ \ {\isachardoublequoteopen}build{\isacharunderscore}{\kern0pt}trees{\isacharprime}{\kern0pt}\ bs\ {\isasymomega}\ k\ i\ I\ {\isacharequal}{\kern0pt}\ {\isacharparenleft}{\kern0pt}\isanewline
\ \ \ \ let\ e\ {\isacharequal}{\kern0pt}\ bs{\isacharbang}{\kern0pt}k{\isacharbang}{\kern0pt}i\ in\ {\isacharparenleft}{\kern0pt}\isanewline
\ \ \ \ case\ pointer\ e\ of\isanewline
\ \ \ \ \ \ Null\ {\isasymRightarrow}\ Some\ {\isacharparenleft}{\kern0pt}{\isacharbrackleft}{\kern0pt}FBranch\ {\isacharparenleft}{\kern0pt}item{\isacharunderscore}{\kern0pt}rule{\isacharunderscore}{\kern0pt}head\ {\isacharparenleft}{\kern0pt}item\ e{\isacharparenright}{\kern0pt}{\isacharparenright}{\kern0pt}\ {\isacharbrackleft}{\kern0pt}{\isacharbrackright}{\kern0pt}{\isacharbrackright}{\kern0pt}{\isacharparenright}{\kern0pt}\isanewline
\ \ \ \ {\isacharbar}{\kern0pt}\ Pre\ pre\ {\isasymRightarrow}\ {\isacharparenleft}{\kern0pt}\isanewline
\ \ \ \ \ \ \ \ do\ {\isacharbraceleft}{\kern0pt}\isanewline
\ \ \ \ \ \ \ \ \ \ pres\ {\isasymleftarrow}\ build{\isacharunderscore}{\kern0pt}trees{\isacharprime}{\kern0pt}\ bs\ {\isasymomega}\ {\isacharparenleft}{\kern0pt}k{\isacharminus}{\kern0pt}{\isadigit{1}}{\isacharparenright}{\kern0pt}\ pre\ {\isacharbraceleft}{\kern0pt}pre{\isacharbraceright}{\kern0pt}{\isacharsemicolon}{\kern0pt}\isanewline
\ \ \ \ \ \ \ \ \ \ those\ {\isacharparenleft}{\kern0pt}map\ {\isacharparenleft}{\kern0pt}{\isasymlambda}f{\isachardot}{\kern0pt}\isanewline
\ \ \ \ \ \ \ \ \ \ \ \ case\ f\ of\isanewline
\ \ \ \ \ \ \ \ \ \ \ \ \ \ FBranch\ N\ fss\ {\isasymRightarrow}\ Some\ {\isacharparenleft}{\kern0pt}FBranch\ N\ {\isacharparenleft}{\kern0pt}fss\ {\isacharat}{\kern0pt}\ {\isacharbrackleft}{\kern0pt}{\isacharbrackleft}{\kern0pt}FLeaf\ {\isacharparenleft}{\kern0pt}{\isasymomega}{\isacharbang}{\kern0pt}{\isacharparenleft}{\kern0pt}k{\isacharminus}{\kern0pt}{\isadigit{1}}{\isacharparenright}{\kern0pt}{\isacharparenright}{\kern0pt}{\isacharbrackright}{\kern0pt}{\isacharbrackright}{\kern0pt}{\isacharparenright}{\kern0pt}{\isacharparenright}{\kern0pt}\isanewline
\ \ \ \ \ \ \ \ \ \ \ \ {\isacharbar}{\kern0pt}\ {\isacharunderscore}{\kern0pt}\ {\isasymRightarrow}\ None\isanewline
\ \ \ \ \ \ \ \ \ \ {\isacharparenright}{\kern0pt}\ pres{\isacharparenright}{\kern0pt}\isanewline
\ \ \ \ \ \ \ \ {\isacharbraceright}{\kern0pt}{\isacharparenright}{\kern0pt}\isanewline
\ \ \ \ {\isacharbar}{\kern0pt}\ PreRed\ p\ ps\ {\isasymRightarrow}\ {\isacharparenleft}{\kern0pt}\isanewline
\ \ \ \ \ \ \ \ let\ ps{\isacharprime}{\kern0pt}\ {\isacharequal}{\kern0pt}\ filter\ {\isacharparenleft}{\kern0pt}{\isasymlambda}{\isacharparenleft}{\kern0pt}k{\isacharprime}{\kern0pt}{\isacharcomma}{\kern0pt}\ pre{\isacharcomma}{\kern0pt}\ red{\isacharparenright}{\kern0pt}{\isachardot}{\kern0pt}\ red\ {\isasymnotin}\ I{\isacharparenright}{\kern0pt}\ {\isacharparenleft}{\kern0pt}p{\isacharhash}{\kern0pt}ps{\isacharparenright}{\kern0pt}\ in\isanewline
\ \ \ \ \ \ \ \ let\ gs\ {\isacharequal}{\kern0pt}\ group{\isacharunderscore}{\kern0pt}by\ {\isacharparenleft}{\kern0pt}{\isasymlambda}{\isacharparenleft}{\kern0pt}k{\isacharprime}{\kern0pt}{\isacharcomma}{\kern0pt}\ pre{\isacharcomma}{\kern0pt}\ red{\isacharparenright}{\kern0pt}{\isachardot}{\kern0pt}\ {\isacharparenleft}{\kern0pt}k{\isacharprime}{\kern0pt}{\isacharcomma}{\kern0pt}\ pre{\isacharparenright}{\kern0pt}{\isacharparenright}{\kern0pt}\ {\isacharparenleft}{\kern0pt}{\isasymlambda}{\isacharparenleft}{\kern0pt}k{\isacharprime}{\kern0pt}{\isacharcomma}{\kern0pt}\ pre{\isacharcomma}{\kern0pt}\ red{\isacharparenright}{\kern0pt}{\isachardot}{\kern0pt}\ red{\isacharparenright}{\kern0pt}\ ps{\isacharprime}{\kern0pt}\ in\isanewline
\ \ \ \ \ \ \ \ map{\isacharunderscore}{\kern0pt}option\ concat\ {\isacharparenleft}{\kern0pt}those\ {\isacharparenleft}{\kern0pt}map\ {\isacharparenleft}{\kern0pt}{\isasymlambda}{\isacharparenleft}{\kern0pt}{\isacharparenleft}{\kern0pt}k{\isacharprime}{\kern0pt}{\isacharcomma}{\kern0pt}\ pre{\isacharparenright}{\kern0pt}{\isacharcomma}{\kern0pt}\ reds{\isacharparenright}{\kern0pt}{\isachardot}{\kern0pt}\isanewline
\ \ \ \ \ \ \ \ \ \ do\ {\isacharbraceleft}{\kern0pt}\isanewline
\ \ \ \ \ \ \ \ \ \ \ \ pres\ {\isasymleftarrow}\ build{\isacharunderscore}{\kern0pt}trees{\isacharprime}{\kern0pt}\ bs\ {\isasymomega}\ k{\isacharprime}{\kern0pt}\ pre\ {\isacharbraceleft}{\kern0pt}pre{\isacharbraceright}{\kern0pt}{\isacharsemicolon}{\kern0pt}\isanewline
\ \ \ \ \ \ \ \ \ \ \ \ rss\ {\isasymleftarrow}\ those\ {\isacharparenleft}{\kern0pt}map\ {\isacharparenleft}{\kern0pt}{\isasymlambda}red{\isachardot}{\kern0pt}\ build{\isacharunderscore}{\kern0pt}trees{\isacharprime}{\kern0pt}\ bs\ {\isasymomega}\ k\ red\ {\isacharparenleft}{\kern0pt}I\ {\isasymunion}\ {\isacharbraceleft}{\kern0pt}red{\isacharbraceright}{\kern0pt}{\isacharparenright}{\kern0pt}{\isacharparenright}{\kern0pt}\ reds{\isacharparenright}{\kern0pt}{\isacharsemicolon}{\kern0pt}\isanewline
\ \ \ \ \ \ \ \ \ \ \ \ those\ {\isacharparenleft}{\kern0pt}map\ {\isacharparenleft}{\kern0pt}{\isasymlambda}f{\isachardot}{\kern0pt}\isanewline
\ \ \ \ \ \ \ \ \ \ \ \ \ \ case\ f\ of\isanewline
\ \ \ \ \ \ \ \ \ \ \ \ \ \ \ \ FBranch\ N\ fss\ {\isasymRightarrow}\ Some\ {\isacharparenleft}{\kern0pt}FBranch\ N\ {\isacharparenleft}{\kern0pt}fss\ {\isacharat}{\kern0pt}\ {\isacharbrackleft}{\kern0pt}concat\ rss{\isacharbrackright}{\kern0pt}{\isacharparenright}{\kern0pt}{\isacharparenright}{\kern0pt}\isanewline
\ \ \ \ \ \ \ \ \ \ \ \ \ \ {\isacharbar}{\kern0pt}\ {\isacharunderscore}{\kern0pt}\ {\isasymRightarrow}\ None\isanewline
\ \ \ \ \ \ \ \ \ \ \ \ {\isacharparenright}{\kern0pt}\ pres{\isacharparenright}{\kern0pt}\isanewline
\ \ \ \ \ \ \ \ \ \ {\isacharbraceright}{\kern0pt}\isanewline
\ \ \ \ \ \ \ \ {\isacharparenright}{\kern0pt}\ gs{\isacharparenright}{\kern0pt}{\isacharparenright}{\kern0pt}\isanewline
\ \ \ \ \ \ {\isacharparenright}{\kern0pt}\isanewline
\ \ {\isacharparenright}{\kern0pt}{\isacharparenright}{\kern0pt}{\isachardoublequoteclose}\isanewline
\isanewline
\isacommand{definition}\isamarkupfalse%
\ build{\isacharunderscore}{\kern0pt}trees\ {\isacharcolon}{\kern0pt}{\isacharcolon}{\kern0pt}\ {\isachardoublequoteopen}{\isacharprime}{\kern0pt}a\ cfg\ {\isasymRightarrow}\ {\isacharprime}{\kern0pt}a\ sentential\ {\isasymRightarrow}\ {\isacharprime}{\kern0pt}a\ bins\ {\isasymRightarrow}\ {\isacharprime}{\kern0pt}a\ forest\ list\ option{\isachardoublequoteclose}\ \isakeyword{where}\isanewline
\ \ {\isachardoublequoteopen}build{\isacharunderscore}{\kern0pt}trees\ {\isasymG}\ {\isasymomega}\ bs\ {\isasymequiv}\isanewline
\ \ \ \ let\ k\ {\isacharequal}{\kern0pt}\ {\isacharbar}{\kern0pt}bs{\isacharbar}{\kern0pt}\ {\isacharminus}{\kern0pt}\ {\isadigit{1}}\ in\isanewline
\ \ \ \ let\ finished\ {\isacharequal}{\kern0pt}\ filter{\isacharunderscore}{\kern0pt}with{\isacharunderscore}{\kern0pt}index\ {\isacharparenleft}{\kern0pt}{\isasymlambda}x{\isachardot}{\kern0pt}\ is{\isacharunderscore}{\kern0pt}finished\ {\isasymG}\ {\isasymomega}\ x{\isacharparenright}{\kern0pt}\ {\isacharparenleft}{\kern0pt}items\ {\isacharparenleft}{\kern0pt}bs{\isacharbang}{\kern0pt}k{\isacharparenright}{\kern0pt}{\isacharparenright}{\kern0pt}\ in\isanewline
\ \ \ \ map{\isacharunderscore}{\kern0pt}option\ concat\ {\isacharparenleft}{\kern0pt}those\ {\isacharparenleft}{\kern0pt}map\ {\isacharparenleft}{\kern0pt}{\isasymlambda}{\isacharparenleft}{\kern0pt}{\isacharunderscore}{\kern0pt}{\isacharcomma}{\kern0pt}\ i{\isacharparenright}{\kern0pt}{\isachardot}{\kern0pt}\ build{\isacharunderscore}{\kern0pt}trees{\isacharprime}{\kern0pt}\ bs\ {\isasymomega}\ k\ i\ {\isacharbraceleft}{\kern0pt}i{\isacharbraceright}{\kern0pt}{\isacharparenright}{\kern0pt}\ finished{\isacharparenright}{\kern0pt}{\isacharparenright}{\kern0pt}{\isachardoublequoteclose}%
\isadelimdocument
%
\endisadelimdocument
%
\isatagdocument
%
\isamarkupsubsection{Termination%
}
\isamarkuptrue%
%
\endisatagdocument
{\isafolddocument}%
%
\isadelimdocument
%
\endisadelimdocument
\isacommand{fun}\isamarkupfalse%
\ build{\isacharunderscore}{\kern0pt}forest{\isacharprime}{\kern0pt}{\isacharunderscore}{\kern0pt}measure\ {\isacharcolon}{\kern0pt}{\isacharcolon}{\kern0pt}\ {\isachardoublequoteopen}{\isacharparenleft}{\kern0pt}{\isacharprime}{\kern0pt}a\ bins\ {\isasymtimes}\ {\isacharprime}{\kern0pt}a\ sentential\ {\isasymtimes}\ nat\ {\isasymtimes}\ nat\ {\isasymtimes}\ nat\ set{\isacharparenright}{\kern0pt}\ {\isasymRightarrow}\ nat{\isachardoublequoteclose}\ \isakeyword{where}\isanewline
\ \ {\isachardoublequoteopen}build{\isacharunderscore}{\kern0pt}forest{\isacharprime}{\kern0pt}{\isacharunderscore}{\kern0pt}measure\ {\isacharparenleft}{\kern0pt}bs{\isacharcomma}{\kern0pt}\ {\isasymomega}{\isacharcomma}{\kern0pt}\ k{\isacharcomma}{\kern0pt}\ i{\isacharcomma}{\kern0pt}\ I{\isacharparenright}{\kern0pt}\ {\isacharequal}{\kern0pt}\ foldl\ {\isacharparenleft}{\kern0pt}{\isacharplus}{\kern0pt}{\isacharparenright}{\kern0pt}\ {\isadigit{0}}\ {\isacharparenleft}{\kern0pt}map\ length\ {\isacharparenleft}{\kern0pt}take\ {\isacharparenleft}{\kern0pt}k{\isacharplus}{\kern0pt}{\isadigit{1}}{\isacharparenright}{\kern0pt}\ bs{\isacharparenright}{\kern0pt}{\isacharparenright}{\kern0pt}\ {\isacharminus}{\kern0pt}\ card\ I{\isachardoublequoteclose}\isanewline
\isanewline
\isacommand{definition}\isamarkupfalse%
\ wf{\isacharunderscore}{\kern0pt}trees{\isacharunderscore}{\kern0pt}input\ {\isacharcolon}{\kern0pt}{\isacharcolon}{\kern0pt}\ {\isachardoublequoteopen}{\isacharparenleft}{\kern0pt}{\isacharprime}{\kern0pt}a\ bins\ {\isasymtimes}\ {\isacharprime}{\kern0pt}a\ sentential\ {\isasymtimes}\ nat\ {\isasymtimes}\ nat\ {\isasymtimes}\ nat\ set{\isacharparenright}{\kern0pt}\ set{\isachardoublequoteclose}\ \isakeyword{where}\isanewline
\ \ {\isachardoublequoteopen}wf{\isacharunderscore}{\kern0pt}trees{\isacharunderscore}{\kern0pt}input\ {\isacharequal}{\kern0pt}\ {\isacharbraceleft}{\kern0pt}\isanewline
\ \ \ \ {\isacharparenleft}{\kern0pt}bs{\isacharcomma}{\kern0pt}\ {\isasymomega}{\isacharcomma}{\kern0pt}\ k{\isacharcomma}{\kern0pt}\ i{\isacharcomma}{\kern0pt}\ I{\isacharparenright}{\kern0pt}\ {\isacharbar}{\kern0pt}\ bs\ {\isasymomega}\ k\ i\ I{\isachardot}{\kern0pt}\isanewline
\ \ \ \ \ \ sound{\isacharunderscore}{\kern0pt}ptrs\ {\isasymomega}\ bs\ {\isasymand}\isanewline
\ \ \ \ \ \ k\ {\isacharless}{\kern0pt}\ {\isacharbar}{\kern0pt}bs{\isacharbar}{\kern0pt}\ {\isasymand}\isanewline
\ \ \ \ \ \ i\ {\isacharless}{\kern0pt}\ {\isacharbar}{\kern0pt}bs{\isacharbang}{\kern0pt}k{\isacharbar}{\kern0pt}\ {\isasymand}\isanewline
\ \ \ \ \ \ I\ {\isasymsubseteq}\ {\isacharbraceleft}{\kern0pt}{\isadigit{0}}{\isachardot}{\kern0pt}{\isachardot}{\kern0pt}{\isacharless}{\kern0pt}{\isacharbar}{\kern0pt}bs{\isacharbang}{\kern0pt}k{\isacharbar}{\kern0pt}{\isacharbraceright}{\kern0pt}\ {\isasymand}\isanewline
\ \ \ \ \ \ i\ {\isasymin}\ I\isanewline
\ \ {\isacharbraceright}{\kern0pt}{\isachardoublequoteclose}\isanewline
\isanewline
\isacommand{lemma}\isamarkupfalse%
\ build{\isacharunderscore}{\kern0pt}trees{\isacharprime}{\kern0pt}{\isacharunderscore}{\kern0pt}termination{\isacharcolon}{\kern0pt}\isanewline
\ \ \isakeyword{assumes}\ {\isachardoublequoteopen}{\isacharparenleft}{\kern0pt}bs{\isacharcomma}{\kern0pt}\ {\isasymomega}{\isacharcomma}{\kern0pt}\ k{\isacharcomma}{\kern0pt}\ i{\isacharcomma}{\kern0pt}\ I{\isacharparenright}{\kern0pt}\ {\isasymin}\ wf{\isacharunderscore}{\kern0pt}trees{\isacharunderscore}{\kern0pt}input{\isachardoublequoteclose}\isanewline
\ \ \isakeyword{shows}\ {\isachardoublequoteopen}{\isasymexists}fs{\isachardot}{\kern0pt}\ build{\isacharunderscore}{\kern0pt}trees{\isacharprime}{\kern0pt}\ bs\ {\isasymomega}\ k\ i\ I\ {\isacharequal}{\kern0pt}\ Some\ fs\ {\isasymand}\ {\isacharparenleft}{\kern0pt}{\isasymforall}f\ {\isasymin}\ set\ fs{\isachardot}{\kern0pt}\ {\isasymexists}N\ fss{\isachardot}{\kern0pt}\ f\ {\isacharequal}{\kern0pt}\ FBranch\ N\ fss{\isacharparenright}{\kern0pt}{\isachardoublequoteclose}%
\isadelimproof
%
\endisadelimproof
%
\isatagproof
%
\endisatagproof
{\isafoldproof}%
%
\isadelimproof
%
\endisadelimproof
%
\begin{isamarkuptext}%
%
\end{isamarkuptext}\isamarkuptrue%
\isacommand{theorem}\isamarkupfalse%
\ termination{\isacharunderscore}{\kern0pt}build{\isacharunderscore}{\kern0pt}tree{\isacharunderscore}{\kern0pt}{\isasymE}arley{\isacharunderscore}{\kern0pt}list{\isacharcolon}{\kern0pt}\isanewline
\ \ \isakeyword{assumes}\ {\isachardoublequoteopen}wf{\isacharunderscore}{\kern0pt}{\isasymG}\ {\isasymG}{\isachardoublequoteclose}\isanewline
\ \ \isakeyword{assumes}\ {\isachardoublequoteopen}nonempty{\isacharunderscore}{\kern0pt}derives\ {\isasymG}{\isachardoublequoteclose}\isanewline
\ \ \isakeyword{assumes}\ {\isachardoublequoteopen}{\isasymG}\ {\isasymturnstile}\ {\isacharbrackleft}{\kern0pt}{\isasymSS}\ {\isasymG}{\isacharbrackright}{\kern0pt}\ {\isasymRightarrow}\isactrlsup {\isacharasterisk}{\kern0pt}\ {\isasymomega}{\isachardoublequoteclose}\isanewline
\ \ \isakeyword{shows}\ {\isachardoublequoteopen}{\isasymexists}fs{\isachardot}{\kern0pt}\ build{\isacharunderscore}{\kern0pt}trees\ {\isasymG}\ {\isasymomega}\ {\isacharparenleft}{\kern0pt}{\isasymE}arley{\isacharunderscore}{\kern0pt}list\ {\isasymG}\ {\isasymomega}{\isacharparenright}{\kern0pt}\ {\isacharequal}{\kern0pt}\ Some\ fs{\isachardoublequoteclose}%
\isadelimproof
%
\endisadelimproof
%
\isatagproof
%
\endisatagproof
{\isafoldproof}%
%
\isadelimproof
%
\endisadelimproof
%
\isadelimdocument
%
\endisadelimdocument
%
\isatagdocument
%
\isamarkupsubsection{Soundness%
}
\isamarkuptrue%
%
\endisatagdocument
{\isafolddocument}%
%
\isadelimdocument
%
\endisadelimdocument
\isacommand{lemma}\isamarkupfalse%
\ wf{\isacharunderscore}{\kern0pt}item{\isacharunderscore}{\kern0pt}tree{\isacharunderscore}{\kern0pt}build{\isacharunderscore}{\kern0pt}trees{\isacharprime}{\kern0pt}{\isacharcolon}{\kern0pt}\isanewline
\ \ \isakeyword{assumes}\ {\isachardoublequoteopen}{\isacharparenleft}{\kern0pt}bs{\isacharcomma}{\kern0pt}\ {\isasymomega}{\isacharcomma}{\kern0pt}\ k{\isacharcomma}{\kern0pt}\ i{\isacharcomma}{\kern0pt}\ I{\isacharparenright}{\kern0pt}\ {\isasymin}\ wf{\isacharunderscore}{\kern0pt}trees{\isacharunderscore}{\kern0pt}input{\isachardoublequoteclose}\isanewline
\ \ \isakeyword{assumes}\ {\isachardoublequoteopen}wf{\isacharunderscore}{\kern0pt}bins\ {\isasymG}\ {\isasymomega}\ bs{\isachardoublequoteclose}\isanewline
\ \ \isakeyword{assumes}\ {\isachardoublequoteopen}k\ {\isacharless}{\kern0pt}\ {\isacharbar}{\kern0pt}bs{\isacharbar}{\kern0pt}{\isachardoublequoteclose}\isanewline
\ \ \isakeyword{assumes}\ {\isachardoublequoteopen}i\ {\isacharless}{\kern0pt}\ {\isacharbar}{\kern0pt}bs{\isacharbang}{\kern0pt}k{\isacharbar}{\kern0pt}{\isachardoublequoteclose}\isanewline
\ \ \isakeyword{assumes}\ {\isachardoublequoteopen}build{\isacharunderscore}{\kern0pt}trees{\isacharprime}{\kern0pt}\ bs\ {\isasymomega}\ k\ i\ I\ {\isacharequal}{\kern0pt}\ Some\ fs{\isachardoublequoteclose}\isanewline
\ \ \isakeyword{assumes}\ {\isachardoublequoteopen}f\ {\isasymin}\ set\ fs{\isachardoublequoteclose}\isanewline
\ \ \isakeyword{assumes}\ {\isachardoublequoteopen}t\ {\isasymin}\ set\ {\isacharparenleft}{\kern0pt}trees\ f{\isacharparenright}{\kern0pt}{\isachardoublequoteclose}\isanewline
\ \ \isakeyword{shows}\ {\isachardoublequoteopen}wf{\isacharunderscore}{\kern0pt}item{\isacharunderscore}{\kern0pt}tree\ {\isasymG}\ {\isacharparenleft}{\kern0pt}item\ {\isacharparenleft}{\kern0pt}bs{\isacharbang}{\kern0pt}k{\isacharbang}{\kern0pt}i{\isacharparenright}{\kern0pt}{\isacharparenright}{\kern0pt}\ t{\isachardoublequoteclose}%
\isadelimproof
%
\endisadelimproof
%
\isatagproof
%
\endisatagproof
{\isafoldproof}%
%
\isadelimproof
%
\endisadelimproof
%
\begin{isamarkuptext}%
%
\end{isamarkuptext}\isamarkuptrue%
\isacommand{lemma}\isamarkupfalse%
\ wf{\isacharunderscore}{\kern0pt}yield{\isacharunderscore}{\kern0pt}tree{\isacharunderscore}{\kern0pt}build{\isacharunderscore}{\kern0pt}trees{\isacharprime}{\kern0pt}{\isacharcolon}{\kern0pt}\isanewline
\ \ \isakeyword{assumes}\ {\isachardoublequoteopen}{\isacharparenleft}{\kern0pt}bs{\isacharcomma}{\kern0pt}\ {\isasymomega}{\isacharcomma}{\kern0pt}\ k{\isacharcomma}{\kern0pt}\ i{\isacharcomma}{\kern0pt}\ I{\isacharparenright}{\kern0pt}\ {\isasymin}\ wf{\isacharunderscore}{\kern0pt}trees{\isacharunderscore}{\kern0pt}input{\isachardoublequoteclose}\isanewline
\ \ \isakeyword{assumes}\ {\isachardoublequoteopen}wf{\isacharunderscore}{\kern0pt}bins\ {\isasymG}\ {\isasymomega}\ bs{\isachardoublequoteclose}\isanewline
\ \ \isakeyword{assumes}\ {\isachardoublequoteopen}k\ {\isacharless}{\kern0pt}\ {\isacharbar}{\kern0pt}bs{\isacharbar}{\kern0pt}{\isachardoublequoteclose}\isanewline
\ \ \isakeyword{assumes}\ {\isachardoublequoteopen}k\ {\isasymle}\ {\isacharbar}{\kern0pt}{\isasymomega}{\isacharbar}{\kern0pt}{\isachardoublequoteclose}\isanewline
\ \ \isakeyword{assumes}\ {\isachardoublequoteopen}i\ {\isacharless}{\kern0pt}\ {\isacharbar}{\kern0pt}bs{\isacharbang}{\kern0pt}k{\isacharbar}{\kern0pt}{\isachardoublequoteclose}\isanewline
\ \ \isakeyword{assumes}\ {\isachardoublequoteopen}build{\isacharunderscore}{\kern0pt}trees{\isacharprime}{\kern0pt}\ bs\ {\isasymomega}\ k\ i\ I\ {\isacharequal}{\kern0pt}\ Some\ fs{\isachardoublequoteclose}\isanewline
\ \ \isakeyword{assumes}\ {\isachardoublequoteopen}f\ {\isasymin}\ set\ fs{\isachardoublequoteclose}\isanewline
\ \ \isakeyword{assumes}\ {\isachardoublequoteopen}t\ {\isasymin}\ set\ {\isacharparenleft}{\kern0pt}trees\ f{\isacharparenright}{\kern0pt}{\isachardoublequoteclose}\isanewline
\ \ \isakeyword{shows}\ {\isachardoublequoteopen}wf{\isacharunderscore}{\kern0pt}yield{\isacharunderscore}{\kern0pt}tree\ {\isasymomega}\ {\isacharparenleft}{\kern0pt}item\ {\isacharparenleft}{\kern0pt}bs{\isacharbang}{\kern0pt}k{\isacharbang}{\kern0pt}i{\isacharparenright}{\kern0pt}{\isacharparenright}{\kern0pt}\ t{\isachardoublequoteclose}%
\isadelimproof
%
\endisadelimproof
%
\isatagproof
%
\endisatagproof
{\isafoldproof}%
%
\isadelimproof
%
\endisadelimproof
%
\begin{isamarkuptext}%
%
\end{isamarkuptext}\isamarkuptrue%
\isacommand{theorem}\isamarkupfalse%
\ wf{\isacharunderscore}{\kern0pt}rule{\isacharunderscore}{\kern0pt}root{\isacharunderscore}{\kern0pt}yield{\isacharunderscore}{\kern0pt}tree{\isacharunderscore}{\kern0pt}build{\isacharunderscore}{\kern0pt}trees{\isacharcolon}{\kern0pt}\isanewline
\ \ \isakeyword{assumes}\ {\isachardoublequoteopen}wf{\isacharunderscore}{\kern0pt}bins\ {\isasymG}\ {\isasymomega}\ bs{\isachardoublequoteclose}\isanewline
\ \ \isakeyword{assumes}\ {\isachardoublequoteopen}sound{\isacharunderscore}{\kern0pt}ptrs\ {\isasymomega}\ bs{\isachardoublequoteclose}\isanewline
\ \ \isakeyword{assumes}\ {\isachardoublequoteopen}{\isacharbar}{\kern0pt}bs{\isacharbar}{\kern0pt}\ {\isacharequal}{\kern0pt}\ {\isacharbar}{\kern0pt}{\isasymomega}{\isacharbar}{\kern0pt}\ {\isacharplus}{\kern0pt}\ {\isadigit{1}}{\isachardoublequoteclose}\isanewline
\ \ \isakeyword{assumes}\ {\isachardoublequoteopen}build{\isacharunderscore}{\kern0pt}trees\ {\isasymG}\ {\isasymomega}\ bs\ {\isacharequal}{\kern0pt}\ Some\ fs{\isachardoublequoteclose}\isanewline
\ \ \isakeyword{assumes}\ {\isachardoublequoteopen}f\ {\isasymin}\ set\ fs{\isachardoublequoteclose}\isanewline
\ \ \isakeyword{assumes}\ {\isachardoublequoteopen}t\ {\isasymin}\ set\ {\isacharparenleft}{\kern0pt}trees\ f{\isacharparenright}{\kern0pt}{\isachardoublequoteclose}\isanewline
\ \ \isakeyword{shows}\ {\isachardoublequoteopen}wf{\isacharunderscore}{\kern0pt}rule{\isacharunderscore}{\kern0pt}tree\ {\isasymG}\ t\ {\isasymand}\ root{\isacharunderscore}{\kern0pt}tree\ t\ {\isacharequal}{\kern0pt}\ {\isasymSS}\ {\isasymG}\ {\isasymand}\ yield{\isacharunderscore}{\kern0pt}tree\ t\ {\isacharequal}{\kern0pt}\ {\isasymomega}{\isachardoublequoteclose}%
\isadelimproof
%
\endisadelimproof
%
\isatagproof
%
\endisatagproof
{\isafoldproof}%
%
\isadelimproof
%
\endisadelimproof
%
\begin{isamarkuptext}%
%
\end{isamarkuptext}\isamarkuptrue%
\isacommand{corollary}\isamarkupfalse%
\ wf{\isacharunderscore}{\kern0pt}rule{\isacharunderscore}{\kern0pt}root{\isacharunderscore}{\kern0pt}yield{\isacharunderscore}{\kern0pt}tree{\isacharunderscore}{\kern0pt}build{\isacharunderscore}{\kern0pt}trees{\isacharunderscore}{\kern0pt}{\isasymE}arley{\isacharunderscore}{\kern0pt}list{\isacharcolon}{\kern0pt}\isanewline
\ \ \isakeyword{assumes}\ {\isachardoublequoteopen}wf{\isacharunderscore}{\kern0pt}{\isasymG}\ {\isasymG}{\isachardoublequoteclose}\isanewline
\ \ \isakeyword{assumes}\ {\isachardoublequoteopen}nonempty{\isacharunderscore}{\kern0pt}derives\ {\isasymG}{\isachardoublequoteclose}\isanewline
\ \ \isakeyword{assumes}\ {\isachardoublequoteopen}build{\isacharunderscore}{\kern0pt}trees\ {\isasymG}\ {\isasymomega}\ {\isacharparenleft}{\kern0pt}{\isasymE}arley{\isacharunderscore}{\kern0pt}list\ {\isasymG}\ {\isasymomega}{\isacharparenright}{\kern0pt}\ {\isacharequal}{\kern0pt}\ Some\ fs{\isachardoublequoteclose}\isanewline
\ \ \isakeyword{assumes}\ {\isachardoublequoteopen}f\ {\isasymin}\ set\ fs{\isachardoublequoteclose}\isanewline
\ \ \isakeyword{assumes}\ {\isachardoublequoteopen}t\ {\isasymin}\ set\ {\isacharparenleft}{\kern0pt}trees\ f{\isacharparenright}{\kern0pt}{\isachardoublequoteclose}\isanewline
\ \ \isakeyword{shows}\ {\isachardoublequoteopen}wf{\isacharunderscore}{\kern0pt}rule{\isacharunderscore}{\kern0pt}tree\ {\isasymG}\ t\ {\isasymand}\ root{\isacharunderscore}{\kern0pt}tree\ t\ {\isacharequal}{\kern0pt}\ {\isasymSS}\ {\isasymG}\ {\isasymand}\ yield{\isacharunderscore}{\kern0pt}tree\ t\ {\isacharequal}{\kern0pt}\ {\isasymomega}{\isachardoublequoteclose}%
\isadelimproof
%
\endisadelimproof
%
\isatagproof
%
\endisatagproof
{\isafoldproof}%
%
\isadelimproof
%
\endisadelimproof
%
\begin{isamarkuptext}%
%
\end{isamarkuptext}\isamarkuptrue%
\isacommand{theorem}\isamarkupfalse%
\ soundness{\isacharunderscore}{\kern0pt}build{\isacharunderscore}{\kern0pt}trees{\isacharunderscore}{\kern0pt}{\isasymE}arley{\isacharunderscore}{\kern0pt}list{\isacharcolon}{\kern0pt}\isanewline
\ \ \isakeyword{assumes}\ {\isachardoublequoteopen}wf{\isacharunderscore}{\kern0pt}{\isasymG}\ {\isasymG}{\isachardoublequoteclose}\isanewline
\ \ \isakeyword{assumes}\ {\isachardoublequoteopen}is{\isacharunderscore}{\kern0pt}sentence\ {\isasymG}\ {\isasymomega}{\isachardoublequoteclose}\isanewline
\ \ \isakeyword{assumes}\ {\isachardoublequoteopen}nonempty{\isacharunderscore}{\kern0pt}derives\ {\isasymG}{\isachardoublequoteclose}\isanewline
\ \ \isakeyword{assumes}\ {\isachardoublequoteopen}build{\isacharunderscore}{\kern0pt}trees\ {\isasymG}\ {\isasymomega}\ {\isacharparenleft}{\kern0pt}{\isasymE}arley{\isacharunderscore}{\kern0pt}list\ {\isasymG}\ {\isasymomega}{\isacharparenright}{\kern0pt}\ {\isacharequal}{\kern0pt}\ Some\ fs{\isachardoublequoteclose}\isanewline
\ \ \isakeyword{assumes}\ {\isachardoublequoteopen}f\ {\isasymin}\ set\ fs{\isachardoublequoteclose}\isanewline
\ \ \isakeyword{assumes}\ {\isachardoublequoteopen}t\ {\isasymin}\ set\ {\isacharparenleft}{\kern0pt}trees\ f{\isacharparenright}{\kern0pt}{\isachardoublequoteclose}\isanewline
\ \ \isakeyword{shows}\ {\isachardoublequoteopen}derives\ {\isasymG}\ {\isacharbrackleft}{\kern0pt}{\isasymSS}\ {\isasymG}{\isacharbrackright}{\kern0pt}\ {\isasymomega}{\isachardoublequoteclose}%
\isadelimproof
%
\endisadelimproof
%
\isatagproof
%
\endisatagproof
{\isafoldproof}%
%
\isadelimproof
%
\endisadelimproof
%
\begin{isamarkuptext}%
%
\end{isamarkuptext}\isamarkuptrue%
%
\isadelimdocument
%
\endisadelimdocument
%
\isatagdocument
%
\isamarkupsubsection{A Word on Performance and Completeness%
}
\isamarkuptrue%
%
\endisatagdocument
{\isafolddocument}%
%
\isadelimdocument
%
\endisadelimdocument
%
\begin{isamarkuptext}%
SNIPPET:

A shared packed parse forest SPPF is a representation designed to reduce the space required to represent multiple derivation
trees for an ambiguous sentence. In an SPPF, nodes which have the same tree below them are shared and nodes which correspond
to different derivations of the same substring from the same non-terminal are combined by creating a packed node for each
family of children. Nodes can be packed only if their yields correspond to the same portion of the input string. Thus, to make it easier
to determine whether two alternates can be packed under a given node, SPPF nodes are labelled with a triple (x,i,j) where
$a_{j+1} \dots a_i$ is a substring matched by x. To obtain a cubic algorithm we use binarised SPPFs which contain intermediate additional
nodes but which are of worst case cubic size. (EXAMPlE SPPF running example???)

We can turn earley's algorithm into a correct parser by adding pointers between items rather than instances of non-terminals, and labelling th epointers
in a way which allows a binariesd SPPF to be constructed by walking the resulting structure. However, inorder to
construct a binarised SPPF we also have to introduce additional nodes for grammar rules of length greater than two,
complicating the final algorithm.%
\end{isamarkuptext}\isamarkuptrue%
%
\isadelimtheory
%
\endisadelimtheory
%
\isatagtheory
%
\endisatagtheory
{\isafoldtheory}%
%
\isadelimtheory
%
\endisadelimtheory
%
\end{isabellebody}%
\endinput
%:%file=05_Earley_Parser.tex%:%
%:%24=9%:%
%:%28=11%:%
%:%38=13%:%
%:%39=13%:%
%:%40=14%:%
%:%41=15%:%
%:%42=16%:%
%:%43=17%:%
%:%44=17%:%
%:%45=18%:%
%:%46=19%:%
%:%47=20%:%
%:%48=21%:%
%:%49=21%:%
%:%50=22%:%
%:%51=23%:%
%:%52=24%:%
%:%53=25%:%
%:%54=25%:%
%:%55=26%:%
%:%56=27%:%
%:%58=29%:%
%:%59=30%:%
%:%60=31%:%
%:%61=31%:%
%:%62=32%:%
%:%63=33%:%
%:%66=36%:%
%:%67=37%:%
%:%68=38%:%
%:%69=38%:%
%:%70=39%:%
%:%77=42%:%
%:%87=44%:%
%:%88=44%:%
%:%89=45%:%
%:%90=46%:%
%:%91=47%:%
%:%92=47%:%
%:%93=48%:%
%:%94=49%:%
%:%95=50%:%
%:%96=50%:%
%:%97=51%:%
%:%100=54%:%
%:%101=55%:%
%:%102=56%:%
%:%103=56%:%
%:%104=57%:%
%:%105=58%:%
%:%106=59%:%
%:%107=59%:%
%:%108=60%:%
%:%111=63%:%
%:%112=64%:%
%:%113=65%:%
%:%114=65%:%
%:%115=66%:%
%:%119=70%:%
%:%120=71%:%
%:%121=72%:%
%:%122=72%:%
%:%123=73%:%
%:%126=76%:%
%:%127=77%:%
%:%128=78%:%
%:%129=78%:%
%:%130=79%:%
%:%131=80%:%
%:%132=81%:%
%:%133=82%:%
%:%134=82%:%
%:%135=83%:%
%:%136=84%:%
%:%137=85%:%
%:%138=86%:%
%:%139=87%:%
%:%140=88%:%
%:%141=89%:%
%:%142=90%:%
%:%143=91%:%
%:%160=98%:%
%:%161=98%:%
%:%162=99%:%
%:%163=100%:%
%:%164=101%:%
%:%165=102%:%
%:%166=103%:%
%:%167=104%:%
%:%168=105%:%
%:%169=106%:%
%:%186=113%:%
%:%187=113%:%
%:%188=113%:%
%:%189=113%:%
%:%190=114%:%
%:%191=115%:%
%:%192=116%:%
%:%193=117%:%
%:%194=118%:%
%:%195=119%:%
%:%212=126%:%
%:%213=126%:%
%:%214=127%:%
%:%215=128%:%
%:%216=129%:%
%:%217=130%:%
%:%218=131%:%
%:%219=132%:%
%:%236=139%:%
%:%237=139%:%
%:%238=140%:%
%:%255=147%:%
%:%256=147%:%
%:%257=148%:%
%:%258=149%:%
%:%259=150%:%
%:%260=151%:%
%:%277=158%:%
%:%278=158%:%
%:%279=159%:%
%:%280=160%:%
%:%281=161%:%
%:%301=167%:%
%:%311=169%:%
%:%312=169%:%
%:%313=170%:%
%:%335=192%:%
%:%336=193%:%
%:%337=194%:%
%:%338=194%:%
%:%339=195%:%
%:%350=202%:%
%:%360=204%:%
%:%361=204%:%
%:%362=205%:%
%:%363=206%:%
%:%364=207%:%
%:%365=207%:%
%:%366=208%:%
%:%372=214%:%
%:%373=215%:%
%:%374=216%:%
%:%375=216%:%
%:%376=217%:%
%:%377=218%:%
%:%401=225%:%
%:%411=227%:%
%:%412=227%:%
%:%413=228%:%
%:%414=229%:%
%:%415=230%:%
%:%416=231%:%
%:%417=232%:%
%:%418=233%:%
%:%435=240%:%
%:%436=240%:%
%:%437=241%:%
%:%438=242%:%
%:%439=243%:%
%:%440=244%:%
%:%441=245%:%
%:%442=246%:%
%:%443=247%:%
%:%460=254%:%
%:%461=254%:%
%:%462=255%:%
%:%463=256%:%
%:%464=257%:%
%:%465=258%:%
%:%466=259%:%
%:%467=260%:%
%:%484=267%:%
%:%485=267%:%
%:%486=268%:%
%:%487=269%:%
%:%488=270%:%
%:%489=271%:%
%:%513=278%:%
%:%523=280%:%
%:%524=280%:%
%:%525=281%:%
%:%526=282%:%
%:%527=283%:%
%:%528=284%:%
%:%548=290%:%
%:%558=292%:%
%:%559=292%:%
%:%560=293%:%
%:%561=294%:%
%:%562=295%:%
%:%563=296%:%
%:%564=296%:%
%:%565=297%:%
%:%566=298%:%
%:%567=299%:%
%:%568=300%:%
%:%569=300%:%
%:%570=301%:%
%:%571=302%:%
%:%581=308%:%
%:%591=310%:%
%:%592=310%:%
%:%593=311%:%
%:%594=312%:%
%:%597=315%:%
%:%598=316%:%
%:%599=317%:%
%:%600=317%:%
%:%601=318%:%
%:%602=319%:%
%:%614=331%:%
%:%617=332%:%
%:%618=332%:%
%:%619=333%:%
%:%647=361%:%
%:%648=362%:%
%:%649=363%:%
%:%650=363%:%
%:%651=364%:%
%:%661=370%:%
%:%671=372%:%
%:%672=372%:%
%:%673=373%:%
%:%674=374%:%
%:%675=375%:%
%:%676=375%:%
%:%677=376%:%
%:%684=383%:%
%:%685=384%:%
%:%686=385%:%
%:%687=385%:%
%:%688=386%:%
%:%689=387%:%
%:%706=394%:%
%:%707=394%:%
%:%708=395%:%
%:%709=396%:%
%:%710=397%:%
%:%711=398%:%
%:%731=404%:%
%:%741=406%:%
%:%742=406%:%
%:%743=407%:%
%:%744=408%:%
%:%745=409%:%
%:%746=410%:%
%:%747=411%:%
%:%748=412%:%
%:%749=413%:%
%:%750=414%:%
%:%767=421%:%
%:%768=421%:%
%:%769=422%:%
%:%770=423%:%
%:%771=424%:%
%:%772=425%:%
%:%773=426%:%
%:%774=427%:%
%:%775=428%:%
%:%776=429%:%
%:%777=430%:%
%:%794=437%:%
%:%795=437%:%
%:%796=438%:%
%:%797=439%:%
%:%798=440%:%
%:%799=441%:%
%:%800=442%:%
%:%801=443%:%
%:%802=444%:%
%:%819=451%:%
%:%820=451%:%
%:%821=452%:%
%:%822=453%:%
%:%823=454%:%
%:%824=455%:%
%:%825=456%:%
%:%826=457%:%
%:%843=464%:%
%:%844=464%:%
%:%845=465%:%
%:%846=466%:%
%:%847=467%:%
%:%848=468%:%
%:%849=469%:%
%:%850=470%:%
%:%851=471%:%
%:%875=479%:%
%:%887=482%:%
%:%888=483%:%
%:%889=484%:%
%:%890=485%:%
%:%891=486%:%
%:%892=487%:%
%:%893=488%:%
%:%894=489%:%
%:%895=490%:%
%:%896=491%:%
%:%897=492%:%
%:%898=493%:%
%:%899=494%:%
%:%900=495%:%

%
\begin{isabellebody}%
\setisabellecontext{{\isadigit{0}}{\isadigit{6}}{\isacharunderscore}{\kern0pt}Examples}%
%
\isadelimtheory
%
\endisadelimtheory
%
\isatagtheory
%
\endisatagtheory
{\isafoldtheory}%
%
\isadelimtheory
%
\endisadelimtheory
%
\isadelimdocument
%
\endisadelimdocument
%
\isatagdocument
%
\isamarkupchapter{The Example in Isabelle%
}
\isamarkuptrue%
%
\endisatagdocument
{\isafolddocument}%
%
\isadelimdocument
%
\endisadelimdocument
\isacommand{definition}\isamarkupfalse%
\ {\isasymepsilon}{\isacharunderscore}{\kern0pt}free\ {\isacharcolon}{\kern0pt}{\isacharcolon}{\kern0pt}\ {\isachardoublequoteopen}{\isacharprime}{\kern0pt}a\ cfg\ {\isasymRightarrow}\ bool{\isachardoublequoteclose}\ \isakeyword{where}\isanewline
\ \ {\isachardoublequoteopen}{\isasymepsilon}{\isacharunderscore}{\kern0pt}free\ {\isasymG}\ {\isasymlongleftrightarrow}\ {\isacharparenleft}{\kern0pt}{\isasymforall}r\ {\isasymin}\ set\ {\isacharparenleft}{\kern0pt}{\isasymRR}\ {\isasymG}{\isacharparenright}{\kern0pt}{\isachardot}{\kern0pt}\ rule{\isacharunderscore}{\kern0pt}body\ r\ {\isasymnoteq}\ {\isacharbrackleft}{\kern0pt}{\isacharbrackright}{\kern0pt}{\isacharparenright}{\kern0pt}{\isachardoublequoteclose}\isanewline
\isanewline
\isacommand{lemma}\isamarkupfalse%
\ {\isasymepsilon}{\isacharunderscore}{\kern0pt}free{\isacharunderscore}{\kern0pt}impl{\isacharunderscore}{\kern0pt}non{\isacharunderscore}{\kern0pt}empty{\isacharunderscore}{\kern0pt}deriv{\isacharcolon}{\kern0pt}\isanewline
\ \ {\isachardoublequoteopen}{\isasymepsilon}{\isacharunderscore}{\kern0pt}free\ {\isasymG}\ {\isasymLongrightarrow}\ N\ {\isasymin}\ set\ {\isacharparenleft}{\kern0pt}{\isasymNN}\ {\isasymG}{\isacharparenright}{\kern0pt}\ {\isasymLongrightarrow}\ {\isasymnot}\ {\isacharparenleft}{\kern0pt}{\isasymG}\ {\isasymturnstile}\ {\isacharbrackleft}{\kern0pt}N{\isacharbrackright}{\kern0pt}\ {\isasymRightarrow}\isactrlsup {\isacharasterisk}{\kern0pt}\ {\isacharbrackleft}{\kern0pt}{\isacharbrackright}{\kern0pt}{\isacharparenright}{\kern0pt}{\isachardoublequoteclose}%
\isadelimproof
%
\endisadelimproof
%
\isatagproof
%
\endisatagproof
{\isafoldproof}%
%
\isadelimproof
\isanewline
%
\endisadelimproof
\isacommand{datatype}\isamarkupfalse%
\ t\ {\isacharequal}{\kern0pt}\ x\ {\isacharbar}{\kern0pt}\ plus\isanewline
\isacommand{datatype}\isamarkupfalse%
\ n\ {\isacharequal}{\kern0pt}\ S\isanewline
\isacommand{datatype}\isamarkupfalse%
\ s\ {\isacharequal}{\kern0pt}\ Terminal\ t\ {\isacharbar}{\kern0pt}\ Nonterminal\ n\isanewline
\isanewline
\isacommand{definition}\isamarkupfalse%
\ nonterminals\ {\isacharcolon}{\kern0pt}{\isacharcolon}{\kern0pt}\ {\isachardoublequoteopen}s\ list{\isachardoublequoteclose}\ \isakeyword{where}\isanewline
\ \ {\isachardoublequoteopen}nonterminals\ {\isacharequal}{\kern0pt}\ {\isacharbrackleft}{\kern0pt}Nonterminal\ S{\isacharbrackright}{\kern0pt}{\isachardoublequoteclose}\isanewline
\isanewline
\isacommand{definition}\isamarkupfalse%
\ terminals\ {\isacharcolon}{\kern0pt}{\isacharcolon}{\kern0pt}\ {\isachardoublequoteopen}s\ list{\isachardoublequoteclose}\ \isakeyword{where}\isanewline
\ \ {\isachardoublequoteopen}terminals\ {\isacharequal}{\kern0pt}\ {\isacharbrackleft}{\kern0pt}Terminal\ x{\isacharcomma}{\kern0pt}\ Terminal\ plus{\isacharbrackright}{\kern0pt}{\isachardoublequoteclose}\isanewline
\isanewline
\isacommand{definition}\isamarkupfalse%
\ rules\ {\isacharcolon}{\kern0pt}{\isacharcolon}{\kern0pt}\ {\isachardoublequoteopen}s\ rule\ list{\isachardoublequoteclose}\ \isakeyword{where}\isanewline
\ \ {\isachardoublequoteopen}rules\ {\isacharequal}{\kern0pt}\ {\isacharbrackleft}{\kern0pt}\isanewline
\ \ \ \ {\isacharparenleft}{\kern0pt}Nonterminal\ S{\isacharcomma}{\kern0pt}\ {\isacharbrackleft}{\kern0pt}Terminal\ x{\isacharbrackright}{\kern0pt}{\isacharparenright}{\kern0pt}{\isacharcomma}{\kern0pt}\isanewline
\ \ \ \ {\isacharparenleft}{\kern0pt}Nonterminal\ S{\isacharcomma}{\kern0pt}\ {\isacharbrackleft}{\kern0pt}Nonterminal\ S{\isacharcomma}{\kern0pt}\ Terminal\ plus{\isacharcomma}{\kern0pt}\ Nonterminal\ S{\isacharbrackright}{\kern0pt}{\isacharparenright}{\kern0pt}\isanewline
\ \ {\isacharbrackright}{\kern0pt}{\isachardoublequoteclose}\isanewline
\isanewline
\isacommand{definition}\isamarkupfalse%
\ start{\isacharunderscore}{\kern0pt}symbol\ {\isacharcolon}{\kern0pt}{\isacharcolon}{\kern0pt}\ s\ \isakeyword{where}\isanewline
\ \ {\isachardoublequoteopen}start{\isacharunderscore}{\kern0pt}symbol\ {\isacharequal}{\kern0pt}\ Nonterminal\ S{\isachardoublequoteclose}\isanewline
\isanewline
\isacommand{definition}\isamarkupfalse%
\ {\isasymG}\ {\isacharcolon}{\kern0pt}{\isacharcolon}{\kern0pt}\ {\isachardoublequoteopen}s\ cfg{\isachardoublequoteclose}\ \isakeyword{where}\isanewline
\ \ {\isachardoublequoteopen}{\isasymG}\ {\isacharequal}{\kern0pt}\ CFG\ nonterminals\ terminals\ rules\ start{\isacharunderscore}{\kern0pt}symbol{\isachardoublequoteclose}\isanewline
\isanewline
\isacommand{definition}\isamarkupfalse%
\ {\isasymomega}\ {\isacharcolon}{\kern0pt}{\isacharcolon}{\kern0pt}\ {\isachardoublequoteopen}s\ list{\isachardoublequoteclose}\ \isakeyword{where}\isanewline
\ \ {\isachardoublequoteopen}{\isasymomega}\ {\isacharequal}{\kern0pt}\ {\isacharbrackleft}{\kern0pt}Terminal\ x{\isacharcomma}{\kern0pt}\ Terminal\ plus{\isacharcomma}{\kern0pt}\ Terminal\ x{\isacharcomma}{\kern0pt}\ Terminal\ plus{\isacharcomma}{\kern0pt}\ Terminal\ x{\isacharbrackright}{\kern0pt}{\isachardoublequoteclose}\isanewline
\isanewline
\isacommand{lemma}\isamarkupfalse%
\ wf{\isacharunderscore}{\kern0pt}{\isasymG}{\isacharcolon}{\kern0pt}\isanewline
\ \ \isakeyword{shows}\ {\isachardoublequoteopen}wf{\isacharunderscore}{\kern0pt}{\isasymG}\ {\isasymG}{\isachardoublequoteclose}%
\isadelimproof
%
\endisadelimproof
%
\isatagproof
%
\endisatagproof
{\isafoldproof}%
%
\isadelimproof
%
\endisadelimproof
%
\begin{isamarkuptext}%
%
\end{isamarkuptext}\isamarkuptrue%
\isacommand{lemma}\isamarkupfalse%
\ is{\isacharunderscore}{\kern0pt}sentence{\isacharunderscore}{\kern0pt}{\isasymomega}{\isacharcolon}{\kern0pt}\isanewline
\ \ \isakeyword{shows}\ {\isachardoublequoteopen}is{\isacharunderscore}{\kern0pt}sentence\ {\isasymG}\ {\isasymomega}{\isachardoublequoteclose}%
\isadelimproof
%
\endisadelimproof
%
\isatagproof
%
\endisatagproof
{\isafoldproof}%
%
\isadelimproof
%
\endisadelimproof
%
\begin{isamarkuptext}%
%
\end{isamarkuptext}\isamarkuptrue%
\isacommand{lemma}\isamarkupfalse%
\ nonempty{\isacharunderscore}{\kern0pt}derives{\isacharcolon}{\kern0pt}\isanewline
\ \ \isakeyword{shows}\ {\isachardoublequoteopen}nonempty{\isacharunderscore}{\kern0pt}derives\ {\isasymG}{\isachardoublequoteclose}%
\isadelimproof
%
\endisadelimproof
%
\isatagproof
%
\endisatagproof
{\isafoldproof}%
%
\isadelimproof
%
\endisadelimproof
%
\begin{isamarkuptext}%
%
\end{isamarkuptext}\isamarkuptrue%
\isacommand{lemma}\isamarkupfalse%
\ correctness{\isacharcolon}{\kern0pt}\isanewline
\ \ \isakeyword{shows}\ {\isachardoublequoteopen}recognizing{\isacharunderscore}{\kern0pt}list\ {\isacharparenleft}{\kern0pt}{\isasymE}arley{\isacharunderscore}{\kern0pt}list\ {\isasymG}\ {\isasymomega}{\isacharparenright}{\kern0pt}\ {\isasymG}\ {\isasymomega}\ {\isasymlongleftrightarrow}\ {\isasymG}\ {\isasymturnstile}\ {\isacharbrackleft}{\kern0pt}{\isasymSS}\ {\isasymG}{\isacharbrackright}{\kern0pt}\ {\isasymRightarrow}\isactrlsup {\isacharasterisk}{\kern0pt}\ {\isasymomega}{\isachardoublequoteclose}%
\isadelimproof
%
\endisadelimproof
%
\isatagproof
%
\endisatagproof
{\isafoldproof}%
%
\isadelimproof
%
\endisadelimproof
%
\isadelimtheory
%
\endisadelimtheory
%
\isatagtheory
%
\endisatagtheory
{\isafoldtheory}%
%
\isadelimtheory
%
\endisadelimtheory
%
\end{isabellebody}%
\endinput
%:%file=06_Examples.tex%:%
%:%24=8%:%
%:%34=10%:%
%:%35=10%:%
%:%36=11%:%
%:%37=12%:%
%:%38=13%:%
%:%39=13%:%
%:%40=14%:%
%:%51=18%:%
%:%54=19%:%
%:%55=19%:%
%:%56=20%:%
%:%57=20%:%
%:%58=21%:%
%:%59=21%:%
%:%60=22%:%
%:%61=23%:%
%:%62=23%:%
%:%63=24%:%
%:%64=25%:%
%:%65=26%:%
%:%66=26%:%
%:%67=27%:%
%:%68=28%:%
%:%69=29%:%
%:%70=29%:%
%:%71=30%:%
%:%74=33%:%
%:%75=34%:%
%:%76=35%:%
%:%77=35%:%
%:%78=36%:%
%:%79=37%:%
%:%80=38%:%
%:%81=38%:%
%:%82=39%:%
%:%83=40%:%
%:%84=41%:%
%:%85=41%:%
%:%86=42%:%
%:%87=43%:%
%:%88=44%:%
%:%89=44%:%
%:%90=45%:%
%:%107=52%:%
%:%108=52%:%
%:%109=53%:%
%:%126=60%:%
%:%127=60%:%
%:%128=61%:%
%:%145=68%:%
%:%146=68%:%
%:%147=69%:%

%
\begin{isabellebody}%
\setisabellecontext{{\isadigit{0}}{\isadigit{7}}{\isacharunderscore}{\kern0pt}Conclusion}%
%
\isadelimtheory
%
\endisadelimtheory
%
\isatagtheory
%
\endisatagtheory
{\isafoldtheory}%
%
\isadelimtheory
%
\endisadelimtheory
%
\isadelimdocument
%
\endisadelimdocument
%
\isatagdocument
%
\isamarkupchapter{Conclusion%
}
\isamarkuptrue%
%
\isamarkupsection{Summary%
}
\isamarkuptrue%
%
\endisatagdocument
{\isafolddocument}%
%
\isadelimdocument
%
\endisadelimdocument
%
\begin{isamarkuptext}%
We formalized and verified an Earley Parser. Our approach and proof development is incremental and
tries to separate the concerns of recognizing and parsing as much as possible.

In an initial step, we introduced an abstract set-based implementation of an Earley recognizer. Our contributions
build upon the foundation of the work on Local Lexing of Obua \cite{Obua:2017} \cite{LocalLexing-AFP},
who formalized the basic building blocks about context-free grammars and derivations. We then proved
the core theorems soundness, completeness and finiteness putting the main focus on the completeness proof
that is mostly inspired by the work of Jones \cite{Jones:1972}.

Subsequently, we refined the Earley recognizer algorithm to a functional executable implementation modeled
after the original imperative implementation of Earley \cite{Earley:1970}. Then we followed once more
Jones footsteps, by not explicitly reproving correctness of the algorithm, but instead replacing set-based
specifications by list-based implementations and proving subsumption in both directions. In the process,
we encountered the hurdles that concern Earley recognizers and grammars containing non-empty derivations, and
ultimately followed Jones one final time and restrained the grammar. Finally, we provided an informal argument
for the running time of our functional implementation and discussed alternative implementation approaches.

We moved on to extend the recognizer into a fully-fledged parser. First, we enriched the recognizer with
the necessary information to construct parse trees in the form of pointers that we modeled after the work
of Scott \cite{Scott:2008} and thus avoided erroneous derivations that occurred in Earley's original
version. We then developed a functional algorithm that utilizes the pointers to construct a single parse tree
and proved its correctness. In a last step, we attempted to generalize our algorithm from a parse tree to
a parse forest. We succeeded in proving soundness, but ultimately abandoned the approach.
Nonetheless, we learned valuable lessons that lay the foundation for future work.

We rounded off the development and highlighted the main theorems by implementing the running example
in Isabelle.%
\end{isamarkuptext}\isamarkuptrue%
%
\isadelimdocument
%
\endisadelimdocument
%
\isatagdocument
%
\isamarkupsection{Future Work%
}
\isamarkuptrue%
%
\endisatagdocument
{\isafolddocument}%
%
\isadelimdocument
%
\endisadelimdocument
%
\begin{isamarkuptext}%
The work on the formalization of Earley parsing is by no means completed. In the following
we sketch - in our opinion - worthwhile future work, roughly presented in decreasing order of importance,
although this might by subject to personal preference.

As mentioned in Chapter \ref{chapter:3}, we followed an operational approach to specify an abstract
set-based implementation of an Earley recognizer. In hindsight, it would have been more elegant to not
follow the approach of Obua but instead define the set of Earley items inductively in Isabelle. This
would probably also simplify some of the subsumption proofs of Chapter \ref{chap:04} that revolve around
function iteration.

From a practical point of view, the last missing piece of this formalization is the construction of a
complete parse forest that represent all possible derivation trees compactly. As discussed in Chapter \ref{chap:05},
the most promising approaches are the algorithms presented by Scott \cite{Scott:2008}, in particular
the algorithm that intertwines the construction of the shared-packed parse forest with the generation of
the Earley bins and thus does not need to concern itself with any technicalities whilst proving termination.
Note that Scott describes her algorithm at a high level of abstraction and consequently any attempted
formalization entails a significant amount of work.

A possible next step are the various opportunities to improve the performance of the algorithm(s).
One non-trivial optimization is another refinement step towards an imperative implementation that incorporates the
necessary performance optimizations to achieve the optimal cubic time and space bounds described in
Chapter \ref{chapter:3}. A formal proof of the complexity bounds should be attempted at this point.
Other performance improvements include incorporating a look-head symbol to prune dead end derivation
trees, and optimizing the representation of the grammar and the bins to enable faster prediction and
completion operations. Or one could reach for the stars and verify the state-of-the-art Earley-based
parsing algorithm Marpa introduced by Kepler \cite{Kegler:2023}. It allows arbitrary context-free
grammars, including grammars with empty derivations, by incorporating the work of Aycock and Horspool \cite{Aycock:2002},
and is based on the algorithms of Earley \cite{Earley:1970} and Leo \cite{Leo:1991} who improves
the running time for right recursions. It claims to run in linear time for nearly all unambiguous grammars,
in particular all LL(k), LR(k), LALR, LRR and operator grammars.

In a last step, one might want to investigate how to incorporate parse tree disambiguation into an
Earley parser. One of strengths of an Earley parser is its ability to work with arbitrary context-free
grammars (at least grammars that contain non-empty derivations in our case). Nonetheless, one might
want to resolve some ambiguities by allowing the user the specify precedence and associativity declarations
restricting the set of allowed parses, as commonly found in parser generators. During our initial
research we investigated possible approaches. The following list is by no means conclusive:
Adams \textit{et al} \cite{Adams:2017} describe a grammar rewriting approach that reinterprets context-free
grammars as tree automata, intersects them with automata encoding the desired restrictions and reinterprets
the results into a context-free grammar. Afroozeh \textit{et al} \cite{Afroozeh:2013} present an approach
of specifying operator precedence based on declarative disambiguation rules basing their implementation
on grammar rewriting. Thorup \cite{Thorup:1996} develops two concrete algorithms for the disambiguation of
grammars based on the idea of excluding a certain set of forbidden sub-parse trees.
Possibly the most interesting approach, since it avoids grammar rewriting and is solely based on parse
tree filtering, is the work of Klint \textit{et al} \cite{Klint:1997}. He propose as framework of filters, that select from
a set of parse trees the intended trees, to describe and compare a wide range of disambiguation problems
in a parser-independent way. Last but not least, Marpa also allows the user to mix classical Chomskyan
parsing and precedence parsing to annotate expressions with precedence and association rules.%
\end{isamarkuptext}\isamarkuptrue%
%
\isadelimtheory
%
\endisadelimtheory
%
\isatagtheory
%
\endisatagtheory
{\isafoldtheory}%
%
\isadelimtheory
%
\endisadelimtheory
%
\end{isabellebody}%
\endinput
%:%file=07_Conclusion.tex%:%
%:%24=8%:%
%:%28=10%:%
%:%40=13%:%
%:%41=14%:%
%:%42=15%:%
%:%43=16%:%
%:%44=17%:%
%:%45=18%:%
%:%46=19%:%
%:%47=20%:%
%:%48=21%:%
%:%49=22%:%
%:%50=23%:%
%:%51=24%:%
%:%52=25%:%
%:%53=26%:%
%:%54=27%:%
%:%55=28%:%
%:%56=29%:%
%:%57=30%:%
%:%58=31%:%
%:%59=32%:%
%:%60=33%:%
%:%61=34%:%
%:%62=35%:%
%:%63=36%:%
%:%64=37%:%
%:%65=38%:%
%:%66=39%:%
%:%75=42%:%
%:%87=45%:%
%:%88=46%:%
%:%89=47%:%
%:%90=48%:%
%:%91=49%:%
%:%92=50%:%
%:%93=51%:%
%:%94=52%:%
%:%95=53%:%
%:%96=54%:%
%:%97=55%:%
%:%98=56%:%
%:%99=57%:%
%:%100=58%:%
%:%101=59%:%
%:%102=60%:%
%:%103=61%:%
%:%104=62%:%
%:%105=63%:%
%:%106=64%:%
%:%107=65%:%
%:%108=66%:%
%:%109=67%:%
%:%110=68%:%
%:%111=69%:%
%:%112=70%:%
%:%113=71%:%
%:%114=72%:%
%:%115=73%:%
%:%116=74%:%
%:%117=75%:%
%:%118=76%:%
%:%119=77%:%
%:%120=78%:%
%:%121=79%:%
%:%122=80%:%
%:%123=81%:%
%:%124=82%:%
%:%125=83%:%
%:%126=84%:%
%:%127=85%:%
%:%128=86%:%
%:%129=87%:%
%:%130=88%:%
%:%131=89%:%
%:%132=90%:%
%:%133=91%:%
%:%134=92%:%



% \appendix{}

%\microtypesetup{protrusion=false}
%\listoffigures{}
%\listoftables{}
%\microtypesetup{protrusion=true}
\printbibliography{}

\end{document}
